\documentclass[11pt]{article}
\pdftrailerid{}
\pdfsuppressptexinfo=15
\pdfgentounicode=1
\pdfglyphtounicode{tfm:lmsy10/mapsto}{007C}
\pdfglyphtounicode{tfm:lmex10/parenleftbig}{0028}
\pdfglyphtounicode{tfm:lmex10/parenrightbig}{0029}
\pdfglyphtounicode{tfm:lmex10/braceleftbig}{007B}
\pdfglyphtounicode{tfm:lmex10/bracerightbig}{007D}
\pdfglyphtounicode{tfm:lmex10/vextendsingle}{23D0}
\pdfglyphtounicode{tfm:lmex10/parenleftBig}{0028}
\pdfglyphtounicode{tfm:lmex10/parenrightBig}{0029}
\pdfglyphtounicode{tfm:lmex10/parenleftbigg}{0028}
\pdfglyphtounicode{tfm:lmex10/parenrightbigg}{0029}
\pdfglyphtounicode{tfm:lmex10/parenleftBigg}{0028}
\pdfglyphtounicode{tfm:lmex10/parenrightBigg}{0029}
\pdfglyphtounicode{tfm:lmex10/bracketleftBigg}{005B}
\pdfglyphtounicode{tfm:lmex10/bracketrightBigg}{005D}
\pdfglyphtounicode{tfm:lmex10/floorleftBigg}{230A}
\pdfglyphtounicode{tfm:lmex10/floorrightBigg}{230B}
\pdfglyphtounicode{tfm:lmex10/parenlefttp}{239B}
\pdfglyphtounicode{tfm:lmex10/parenrighttp}{239E}
\pdfglyphtounicode{tfm:lmex10/bracelefttp}{23A7}
\pdfglyphtounicode{tfm:lmex10/braceleftbt}{23A9}
\pdfglyphtounicode{tfm:lmex10/braceleftmid}{23A8}
\pdfglyphtounicode{tfm:lmex10/parenleftbt}{239D}
\pdfglyphtounicode{tfm:lmex10/parenrightbt}{23A0}
\pdfglyphtounicode{tfm:lmex10/intersectiontext}{22C2}
\pdfglyphtounicode{tfm:lmex10/summationdisplay}{2211}
\pdfglyphtounicode{tfm:lmex10/uniondisplay}{22C3}
\pdfglyphtounicode{tfm:lmex10/intersectiondisplay}{22C2}
\usepackage[T1]{fontenc}
\usepackage{lmodern}
\usepackage[margin=1in]{geometry}
\usepackage{amsmath,amssymb,amsthm,mathtools}
\usepackage{enumitem,booktabs,array}
\usepackage{xcolor,graphicx,tikz}
\usetikzlibrary{arrows.meta,positioning,calc,decorations.pathreplacing}
\usepackage{microtype}
\usepackage[numbers,sort&compress]{natbib}
\usepackage{xurl}
\usepackage[colorlinks=true,linkcolor=blue!45!black,citecolor=blue!45!black,urlcolor=blue!45!black]{hyperref}
\usepackage[figure]{hypcap}
\let\originalbibsection\bibsection
\renewcommand{\bibsection}{%
  \originalbibsection
  \addcontentsline{toc}{section}{\refname}}
\newcommand{\F}{\mathbf F}
\newcommand{\eps}{\epsilon}
\DeclareMathOperator{\sht}{ht}
\newtheorem{theorem}{Theorem}[section]
\newtheorem{lemma}[theorem]{Lemma}
\newtheorem{proposition}[theorem]{Proposition}
\newtheorem{corollary}[theorem]{Corollary}
\theoremstyle{definition}
\newtheorem{definition}[theorem]{Definition}
\theoremstyle{remark}
\newtheorem{remark}[theorem]{Remark}
\numberwithin{equation}{section}
\setlist[enumerate]{itemsep=2pt,topsep=5pt}
\setlist[itemize]{itemsep=2pt,topsep=5pt}
\setlength{\emergencystretch}{1em}


\title{Generalized Star Height at Most Four}
\author{OpenAI}
\date{September 25, 2026}
\hypersetup{pdftitle={Generalized Star Height at Most Four},pdfauthor={OpenAI}}
\begin{document}
\maketitle
\begin{abstract}
Every regular language over a finite alphabet has generalized star height
at most four over that same alphabet. We give a complete construction using
an affine correction that hides one interval product, a finite clock, and
several scales for moving boundaries through periodic words.
\end{abstract}
\section{Introduction}\label{sec:introduction}

How deeply must Kleene stars be nested to describe a regular language when
complement is also available? The answer depends on the permitted
operations. In this paper an expression over a fixed finite alphabet
$\Sigma$ starts from $0$, $1$, and the letters of $\Sigma$, denoting
$\varnothing$, $\{\eps\}$, and the corresponding singleton languages.
It uses union, concatenation, complement in $\Sigma^*$, and Kleene star.
Its height is given by
\[
\begin{aligned}
 \sht(0)=\sht(1)=\sht(a)&=0,\\
 \sht(P\cup Q)=\sht(PQ)&=\max\{\sht(P),\sht(Q)\},\\
 \sht(\neg P)&=\sht(P),&
 \sht(P^*)&=1+\sht(P).
\end{aligned}
\]
The \emph{generalized star height} $h_\Sigma(L)$ is the minimum height of
an expression defining $L\subseteq\Sigma^*$. Finite Boolean operations
preserve any common height bound. Finite languages have height zero, as
does $\Sigma^*=\neg0$. The complement operation keeps the height of its
operand; it does not remove stars already nested in that operand.

\begin{theorem}\label{thm:main}
For every finite alphabet $\Sigma$ and every regular language
$L\subseteq\Sigma^*$, one has $h_\Sigma(L)\le4$.
\end{theorem}

We prove the theorem by a complete construction over the given alphabet.
The bound is independent of $\Sigma$, $L$, and the number of states in a
recognizing automaton. Our finite choices can be found by exhaustive
search; we assert no useful bound on their sizes or on expression length.

\subsection{The problem and the three constructions}

The ordinary problem arose from expressions without complement. Eggan
related their star height to transition graphs~\cite{Eggan1963};
Dejean and Sch\"utzenberger established unbounded ordinary star height over
a binary alphabet~\cite{DejeanSchutzenberger1966}. Those lower bounds do
not carry over when complement is added. In the Boolean model,
Sch\"utzenberger characterized the height-zero, or star-free, languages
as exactly those recognized by aperiodic finite
monoids~\cite{Schutzenberger1965}.

Height-one results extend well beyond the aperiodic case. Henneman proved
that every language recognized by a finite commutative group has height
at most one, as recalled in \cite[Theorem~3.4]{PST1992}. Pin, Straubing,
and Th\'erien extended this conclusion to finite nilpotent groups of class
two~\cite[Theorem~7.3]{PST1992}. These algebraic results do not supply a
uniform bound for arbitrary finite monoids. Absolute boundedness and the
assertion that height one always suffices are different questions;
Straubing explicitly formulates both~\cite[p.~537]{Straubing2002}.
Theorem~\ref{thm:main} supplies an absolute bound. It does not exhibit a
language of height greater than one or decide whether one always suffices.

The alphabet convention is part of this problem. An alphabetic morphism
sends each letter to a letter or the empty word. Pin, Straubing, and
Th\'erien prove that inverse alphabetic morphisms do not increase generalized
height~\cite[Corollary~4.7]{PST1992}. They also express every regular
language as the inverse image, under a general morphism, of a language
of ordinary height at most one~\cite[Theorem~8.1]{PST1992}. Preservation of every
generalized-height bound under arbitrary inverse morphisms would therefore
already imply the bound one. Our component expressions are instead constructed
directly over the given alphabet; finite tags index tests and never become
new input letters.

This paper and two companions give different ways to bound the nesting.
\emph{Finite Monoid Computations and a Uniform Generalized Star-Height
Bound}~\cite[Theorem~1.1]{Uniform13} gives height thirteen using prediction
trials and a computation that records the full history of earlier tests.
The present construction gives height four using an affine correction and
several scales at which product boundaries may move. The companion
\emph{Generalized Star Height at Most Three}~\cite[Theorem~1.1]{Height3}
gives the stronger numerical bound three through arbitrary-function shift
tables and a decoder that controls its entire set of candidate cuts.
Each paper proves its own required algebra, timing statements, and
original-alphabet compilation. The purpose here is to develop the complete
multiscale method, including its two different residual decoders.

The split and affine-recovery framework is also developed in
\cite[Sections~2--4]{Uniform13}. Our local split lemma permits both factors
to have height at most one, and we prove that precise form below. The
finite probability argument behind our affine correction is an application
of the dependency-graph method of Erd\H{o}s and Lov\'asz
\cite[Section~2]{ErdosLovasz1975}. Our finite-window markers use the
marker-versus-periodicity principle associated with Krieger's marker lemma
\cite[Lemma~2]{Krieger1982}; the particular local rule and every periodic
transport claim needed here are proved in full. These comparisons concern
methods and mathematical scope: neither companion's numerical theorem is
a premise of this proof.

\subsection{From products to three splits}

A deterministic automaton turns each word into a transformation of a finite
set. These transformations form a finite monoid $M$, and their product in
reading order determines acceptance. We therefore seek low-height tests
for the value of a morphism $T:\Sigma^*\to M$. Cancellation and
commutativity are not assumed.

The word will be parsed into nonempty pieces called \emph{episodes},
followed by a remainder with no episode prefix. Each episode has a unique
end, determined by a bounded marker search or by the first departure from
a periodic continuation. A graph test records the input and output of a
fixed finite-state update on a sequence of episodes. To construct such a
test, we split selected episodes into a prefix test and a suffix test,
and use an already available graph test to skip the other episodes. The
two factors choose their finite information independently. Soundness must
therefore hold for every accepted pair, including pairs with different
guesses and a proposed suffix that ends too soon or too late.

The first difficulty is the product of the interval containing the cut:
neither side can read it in full. Suppose a segment has interval products
$(d_1,\ldots,d_l)\in M^l$ and carries a state in a finite vector space.
We choose one correction function $\beta$ on these tuples so that, for
every input state and every tuple, the corrected output is constant when
one coordinate $d_j$ varies. The prefix can then transmit the input state
and the earlier products; the suffix tests constancy while supplying the
later products. The unknown product $d_j$ need not be transmitted. Crucially,
the same $\beta$ works simultaneously for all inputs and tuples.

The second difficulty is timing. A suffix knows where it begins but may
not know the origin of the episode that produced it. A finite clock
restricts the possible pairs of prefix and suffix information to matching
tags or one uniquely determined mismatch. A deterministic affine
correction handles that mismatch. For matching tags, we choose boundaries
that remain fixed near markers and can move through a periodic region
without changing ordered monoid products. Several scales of movement and
a polynomial-indexed set of cut positions make the number of ambiguous
origin pairs smaller than the number of copies of every needed cut.

There are two remaining timing obstructions: failure of the clock or of
all inner scales, and instability of the later end search. Two buffers
handle them by different decoders. One first forces a common exact end and
then decodes an origin; the other first decodes an origin using bounded
data and then checks that origin's exact end. Their physical order lets
each suffix read the complete later update table. Those tables are defined
at every hypothetical monoid value, including impossible prefix values.
Lifting them to linear maps on a larger basis permits recovery on whole
episode sequences.

These operations lead to three applications of one split identity. Starting
from the empty skipped class, their heights are $0\to2\to3\to4$; the
intervening affine and basis recoveries use only finite Boolean operations.
To justify this count we must also control every component language on its
own, under false guesses as well as true ones. Each has only one unbounded
periodic continuation and hence a finite cover by word templates
$a b^h c$. Ultimately periodic tests on $h$ compile directly over $\Sigma$.
No tag or state is introduced as an extra input letter.

\paragraph{Organization.}
Section~\ref{sec:algebra} gives the monoid reduction, independent-witness
split principle, and affine-line lemma. Section~\ref{sec:clock} constructs
the finite clock and proves its roster-independent bound.
Section~\ref{sec:scales} builds the roster, scales, and moving boundaries
and counts ambiguous origins. Section~\ref{sec:episodes} places all windows
and proves exact end detection. Section~\ref{sec:splits} defines the total
updates and performs the three splits. Section~\ref{sec:compilation}
compiles the individual tests, handles the remainder, and completes the
height calculation.

\section{Algebra and split computations}\label{sec:algebra}

This section separates the algebraic work from the later geometry. We first
represent a word by its finite-monoid value. We then show how a sound split
extends graph tests from a skipped class of episodes to a larger class,
and how affine corrections can be removed from sequence computations.
The final lemma supplies the correction that lets a split omit one interval
product.

We may assume that $\Sigma$ is nonempty: otherwise its only languages
are $\varnothing$ and $\{\eps\}$, both of generalized star height zero.
Fix a morphism
\[
 T:\Sigma^*\longrightarrow M
\]
into a finite monoid, and put $m=|M|$ and $N=m!$.
It suffices to construct bounded-height expressions for the fibers of
$T$. Indeed, a deterministic finite automaton supplies such a morphism
by assigning to a word its transformation of the state set, and its
accepted language is a finite union of fibers.

Letters occupy unit intervals with integer endpoints. For an available
interval $[a,b)$, with $a\le b$, write $T_{a,b}$ for the product of its
letters; in particular $T_{a,a}=1_M$. Products and actions are written
in reading order, so $T_{a,c}=T_{a,b}T_{b,c}$ when $a\le b\le c$.
Let $V=\F_2^M$ have basis $(e_c)_{c\in M}$. Right multiplication in $M$ induces the
linear right action $e_c d=e_{cd}$, extended to all of $V$.
The image $e_{1_M}d=e_d$ distinguishes every $d\in M$.
We will also use other finite-dimensional vector spaces over $\F_2$,
with linear right actions written $v\rho(d)$.

One elementary consequence of finiteness will be used when boundaries
move through periodic words:
\begin{equation}\label{eq:monoid-periodicity}
 c^{h+N}=c^h\qquad(c\in M,\ h\ge m).
\end{equation}
To see this, choose $0\le a<b\le m$ with $c^a=c^b$.
Multiplication by further powers gives period $b-a$ from exponent $a$
onward, and $b-a$ divides $N$. This also covers $m=N=1$.

\subsection{Splitting a sequence of episodes}

A \emph{prefix code} is a set $E\subseteq\Sigma^+$ in which no word is
a proper prefix of another. Every word in $E^*$ has a unique
factorization into elements of $E$: compare the first factors, which
must coincide by the prefix condition, and continue. We call the
elements of $E$ \emph{episodes}.

Suppose each episode $e$ carries a deterministic update $G_e$ of a
fixed finite state set $S$. A \emph{graph language} from $x$ to $y$ on
$D^*$, where $D\subseteq E$, consists of those words in $D^*$ whose
successive episode updates send $x$ to $y$. The empty word carries the
identity update. Throughout, a graph language includes its indicated
domain condition; it accepts no word outside $D^*$.

The next lemma receives two independently defined languages for the factors
of a split. It assumes universal soundness of their concatenations but asks
for a successful split only outside the skipped class. In particular,
completeness does not license either side to trust the other side's guesses.
The common update on an episode must be fixed before either guess is made.

\begin{lemma}[Split principle]\label{lem:split}
Let $D'\subseteq D\subseteq E$, where $E$ is a prefix code, and let
$(G_e)_{e\in D}$ be deterministic updates of a finite set $S$.
For $x,y\in S$, suppose languages $P_x,Q_y$ satisfy the following
conditions.
\begin{enumerate}[label=\textup{(\roman*)}]
 \item If a word in $D^*$ has a prefix $uv$ with $u\in P_x$ and
 $v\in Q_y$, then its first episode $e$ exists, $uv=e$, and
 $G_e(x)=y$.
 \item For every $e\in D\setminus D'$ and every $x\in S$,
 $e\in P_xQ_{G_e(x)}$.
\end{enumerate}
Let $W_{x,y}$ be the graph languages for the same updates on $(D')^*$.
Then the graph language from $x$ to $y$ on $D^*$ is
\begin{equation}\label{eq:split-expression}
\begin{split}
D^*\cap\Bigg[W_{x,y}\ \cup\ &
 \left(\bigcup_{a\in S}W_{x,a}P_a\right)
 \left(\bigcup_{b,c\in S}Q_bW_{b,c}P_c\right)^*\\[-2pt]
 &\hspace{22mm}\cdot
 \left(\bigcup_{f\in S}Q_fW_{f,y}\right)\Bigg].
\end{split}
\end{equation}
If $D$, all $P_x$, and all $Q_y$ have height at most one, and all
$W_{x,y}$ have height at most $b\ge0$, then these new graph languages
have height at most $\max\{2,b+1\}$.
\end{lemma}

\begin{proof}
\emph{Soundness of the expression.}
Consider a factorization accepted by the expression on the right of
\eqref{eq:split-expression}. The outer intersection puts the word in
$D^*$, so it has its true episode boundaries. A $W$-factor starting
at such a boundary consumes a sequence of complete episodes:
its factorization into members of $D'\subseteq E$ agrees with the
ambient factorization by the prefix-code property.

In the second term, pair each $P$-factor with the following $Q$-factor.
Condition~\textup{(i)} says that this pair consumes exactly the next
episode and has the indicated input and output states. Induction
through the alternating $W$-factors and these pairs proves both the
boundary alignment and the required composite update. This includes
the case in which the central star contributes no factors, leaving
one explicit split. The first term $W_{x,y}$ is already correct.

\emph{Completeness and height.}
Conversely, take a word in $D^*$ with the required update. Split each
episode outside $D'$ using condition~\textup{(ii)}, and absorb each
intervening sequence of $D'$-episodes into its graph language $W$.
This gives the second term, unless there are no explicit splits, in
which case it gives $W_{x,y}$. Thus the displayed language is exact.

For the height estimate, $D^*$ has height at most two. The central
star has height at most $1+\max\{1,b\}$, and all other operations are
finite unions, concatenations, or intersections. The stated bound
follows.
\end{proof}

\subsection{Recovering linear parts}

The split identity has reduced the language problem to episode updates.
We are free to add an episode-dependent affine correction if we can later
recover the desired linear action. The next identity does that recovery
on the same complete word in both tests. It is therefore valid after
arbitrarily many episode updates, not merely after one.

\begin{lemma}[Affine recovery]\label{lem:affine-recovery}
Let $U$ be a finite-dimensional vector space over $\F_2$ and let
$D$ be a subset of a prefix code. Suppose the update for $e\in D$ is
\[
 G_e(v)=vL_e+\sigma_e,
\]
where $L_e$ is linear and $\sigma_e\in U$, both determined by the
episode. If $R_{v,z}$ are the graph languages of the affine composite
on $D^*$, then the graph language of the product of the linear parts,
from $v$ to $y$, is
\[
 \bigcup_{f\in U}\bigl(R_{v,y+f}\cap R_{0,f}\bigr).
\]
In particular this recovery does not increase a common height bound.
\end{lemma}

\begin{proof}
Applying $v\mapsto vL+c$ and then $v\mapsto vA+d$ gives
$v\mapsto vLA+cA+d$. Hence the composite on any word of $D^*$ has
the form $v\mapsto vL+c$, where $L$ is the product of the episode
linear parts in reading order. Its output at zero is $c$.
The displayed intersection says that the two outputs are $y+f$ and
$f$, which is equivalent to $vL=y$. The union ranges over the finite
set of possible $f$. Only finite Boolean operations are used.
\end{proof}

\subsection{An affine correction that hides one interval product}

At an interior cut, neither factor can read the product across the
cut. The following finite lemma provides a correction for which one
interval product can be omitted. The omitted interval may depend on
the input state and the full tuple, but the correction itself does
not depend on the input state.

\begin{lemma}[Affine-line constancy]\label{lem:affine-line}
Let $U$ be a finite-dimensional vector space over $\F_2$ and let
$\rho$ be a linear right action of $M$ on $U$. There are an integer
$l\ge1$ and a function $\beta:M^l\to U$ such that, for every
$v\in U$ and every $\mathbf d=(d_1,\ldots,d_l)\in M^l$, some
$j\in\{1,\ldots,l\}$ makes the function
\[
 a\longmapsto
 v\rho(d_1\cdots d_{j-1}a d_{j+1}\cdots d_l)
 +\beta(d_1,\ldots,d_{j-1},a,d_{j+1},\ldots,d_l)
\]
constant on $M$.
\end{lemma}

\begin{proof}
We first bound the probability of failure for one state and tuple, then
show that all these failures can be avoided simultaneously. The length
$l$ is chosen only after both estimates are available.

Write $q_U=|U|$. If $m=1$ or $q_U=1$, take $l=1$ and $\beta=0$.
We therefore assume $m,q_U>1$.
For an integer $l$ to be chosen, choose the values of $\beta$
independently and uniformly in $U$ at all points of $M^l$.
For each $(v,\mathbf d)$, let $\mathcal B_{v,\mathbf d}$ be the event
that no coordinate line through $\mathbf d$ has the asserted
constancy.

Condition on $\beta(\mathbf d)$. Constancy along a specified line
then prescribes the values of $\beta$ at its other $m-1$ points, so
has probability $q_U^{1-m}$. Different coordinate lines through
$\mathbf d$ have disjoint sets of noncentral points. Their constancy
events are therefore independent under this conditioning, and the
conditional probability of failure does not depend on the center
value. Consequently
\[
 \Pr(\mathcal B_{v,\mathbf d})
 =p_l:=(1-q_U^{1-m})^l.
\]

An event $\mathcal B_{v,\mathbf d}$ uses only the random values at
points of Hamming distance at most one from $\mathbf d$. It is thus
independent of the joint family of events whose centers have distance
greater than two from $\mathbf d$. The number of potentially
dependent events is at most
\[
 q_U\left(1+l(m-1)+\binom l2(m-1)^2\right)-1
 \ \le\ D_l:=q_U(1+lm+l^2m^2).
\]
This count includes all possible input vectors at each nearby center.

We use a local-lemma argument \citep[Section~2, pp.~616--617]{ErdosLovasz1975},
recalling the required finite probability estimate with its proof.
Suppose a finite family of events has a dependency graph of maximum
degree at most $D\ge1$: each event is independent of the joint family
of its nonneighbors. If every event has probability at most
$u(1-u)^D$, where $0<u<1$, then their simultaneous nonoccurrence
has positive probability.

For completeness, induct on the size of a conditioning set to show
that its nonoccurrence has positive probability and that the
conditional probability of any other event is at most $u$.
The assertion with an empty conditioning set follows from the
probability bound. Positivity for a set of size $k$ follows by
removing one member and applying the conditional bound for size
$k-1$. To prove the conditional bound for an event $A$ and a set
$S$ of size $k$ not containing it, partition $S$ into its neighbors
$S_1$ and nonneighbors $S_0$. If $S_1$ is empty, joint independence
gives the unconditional bound. Otherwise,
\[
 \Pr\left(A\,\middle|\,\bigcap_{B\in S}B^c\right)
 \le
 \frac{\Pr\left(A\,\middle|\,\bigcap_{B\in S_0}B^c\right)}
 {\Pr\left(\bigcap_{B\in S_1}B^c\,\middle|\,
                    \bigcap_{B\in S_0}B^c\right)}
 \le \frac{\Pr(A)}{(1-u)^{|S_1|}}
 \le u.
\]
The denominator is bounded by revealing the members of $S_1$
successively: every conditional set then has size at most $k-1$,
so the induction applies. The numerator uses joint independence
from $S_0$. This completes the induction and proves the estimate.

Apply the estimate with $D=D_l$ and $u=1/(2D_l)$. Bernoulli's
inequality gives
\[
 u(1-u)^{D_l}
 =\frac1{2D_l}\left(1-\frac1{2D_l}\right)^{D_l}
 \ge\frac1{4D_l}.
\]
The quantity $p_l$ decays exponentially in $l$, whereas $D_l$ grows
quadratically. Choosing $l$ large enough makes $p_l\le1/(4D_l)$.
There is therefore a choice of $\beta$ for which no failure event
occurs. Fix such a choice.
\end{proof}

\begin{remark}[Uniformity and finite choice]\label{rem:finite-affine-choice}
The choices of $l$ and $\beta$ precede the state $v$ and tuple $\mathbf d$.
Only the successful coordinate $j$ may depend on them. The proof is also
an effective finite existence argument: increase $l$ until its displayed
probability inequality holds, then enumerate the finitely many functions
$M^l\to U$ and test the stated property for every state and tuple. At least
one function passes. This procedure asserts no useful running-time bound.
\end{remark}

To apply the lemma to a word segment, choose ordered checkpoints
$B_0<\cdots<B_l$ and set $d_j=T_{B_{j-1},B_j}$.
Process the segment by $\rho$, then add $\beta(\mathbf d)$ at $B_l$.
If a cut lies in interval $j$, the prefix can transmit the vector
$v$ at $B_0$ and the preceding products $d_1,\ldots,d_{j-1}$.
The suffix supplies $d_{j+1},\ldots,d_l$ and checks constancy on the
missing-coordinate line. When that check holds, it obtains the
correct vector at $B_l$ without knowing $d_j$. For every actual
tuple and input vector at least one interval permits this procedure.
Because the added vector is a function of the actual tuple alone,
the update of the whole segment is affine and remains eligible for
Lemma~\ref{lem:affine-recovery}.

\section{A finite tag clock}
\label{sec:clock}

The affine-line lemma requires the two factors to communicate finite data.
We call a possible collection of these data a \emph{tag}. This section
constructs a clock that controls independent tag choices by testing the
two length residues, measured from scheduled reference positions. Its relation permits either matching
tags or a single specified mismatch. The latter can later be absorbed into
an affine correction. A second property bounds ambiguous pairs of candidate
origins uniformly in their common translation. That bound must depend only
on the tags, since the number of candidate origins will be chosen afterward.

\begin{lemma}[Finite tag clock]\label{lem:clock}
Let $\Lambda$ be a finite nonempty set, let $\mathcal R$ be any finite set,
and let $Q\geq 1$ be an integer.  Put
\[
 F_{\rm cl}=400,
 \qquad
 C_{\rm cl}=2^{|\Lambda|}\bigl(1+4F_{\rm cl}^{2}\bigr).
\]
There exist a positive integer $n$, all of whose prime factors exceed $Q$,
sets $I_\alpha,J_\alpha\subseteq\mathbb Z/n\mathbb Z$ for
$\alpha\in\Lambda$, a set $W\subseteq\mathbb Z/n\mathbb Z$, and residues
$v_p\in\mathbb Z/n\mathbb Z$ for $p\in\mathcal R$ with the following
properties.
\begin{enumerate}[label=\textup{(\roman*)}]
\item For every $s\in\mathbb Z/n\mathbb Z$, the relation
\[
 \mathcal E(s)=
 \{(\alpha,\alpha')\in\Lambda^2:
             s\in I_\alpha+J_{\alpha'}\}
\]
is either a subset of the diagonal or a singleton consisting of an
off-diagonal pair.  If $s\notin W$, it is the full diagonal.
\item For every $u\in\mathbb Z/n\mathbb Z$,
\[
 \#\{(p,p')\in\mathcal R^2:p\ne p',\quad
                 u+v_{p'}-v_p\in W\}\leq C_{\rm cl}.
\]
\end{enumerate}
In particular, $C_{\rm cl}$ depends only on $|\Lambda|$.
\end{lemma}

\begin{proof}
The construction uses one block of coordinates for every left/right
assignment of the tags. First-coordinate restrictions exclude competing
tag pairs, while a batch of additional coordinates excludes reverse pairs
and controls translated roster differences. We prove these two tasks
separately and only then identify the product group with one cyclic group.

\emph{Coordinate blocks and their sets.}
We first construct the sets in a finite product of cyclic groups.  For each
map $\chi:\Lambda\to\{\mathrm L,\mathrm R\}$, reserve one first
coordinate and a batch of coordinates.  The batch contains a dedicated
coordinate for every ordered list
\[
 \bigl((p_a,p'_a)\bigr)_{a=0}^{F_{\rm cl}-1}
\]
of ordered pairs of distinct elements of $\mathcal R$ such that the
underlying unordered edges are distinct and form a forest.  Add a spare
coordinate if necessary so that the batch has positive odd size.  Thus, even
when no such list exists, the batch consists of one spare coordinate.
There are finitely many coordinates.  Give them distinct prime orders $q$,
each satisfying
\[
 q>\max\{Q,400,2^{|\mathcal R|+1}\}.
\]
There are enough primes: the product-plus-one argument supplies a prime
outside any prescribed finite list.  Write $G$ for the resulting product of
cyclic groups.

Identify a coordinate of order $q$ with the subgroup
$q^{-1}\mathbb Z/\mathbb Z$ of the circle
$\mathbb T=\mathbb R/\mathbb Z$.  Write $\|x\|_{\mathbb T}$ for circular
distance to zero, and call a coordinate value \emph{small} if
$\|x\|_{\mathbb T}\leq 1/100$.  Define the two sets of centers
\[
 C_I=\{0,1,3\}/6,
 \qquad C_J=\{2,4,5\}/6
 \quad\text{in }\mathbb T.
\]
For each tag $\alpha$, impose the following constraints in the block
belonging to $\chi$.
\begin{itemize}
\item If $\chi(\alpha)=\mathrm L$, require the first coordinate of
  $I_\alpha$ to be zero and that of $J_\alpha$ to be nonzero.  In every
  batch coordinate, require the $I_\alpha$ value to have distance at most
  $1/16$ from $C_I$, and the $J_\alpha$ value to have distance at most
  $1/16$ from $C_J$.
\item If $\chi(\alpha)=\mathrm R$, require the first coordinate of
  $I_\alpha$ to be nonzero and that of $J_\alpha$ to be zero.  For each of
  $I_\alpha$ and $J_\alpha$, require a strict majority of the batch
  coordinates to be small.
\end{itemize}
The constraints over all blocks define $I_\alpha,J_\alpha\subseteq G$.
Let $W\subseteq G$ be the union, over all $\chi$, of the following two
obstructions: the first coordinate is zero; or no coordinate of the batch
is small.

\emph{Sums and differences inside one batch.}
We verify the sum and difference properties of a batch.  For a left tag,
every target in each cyclic coordinate is a sum of allowed values.  Indeed,
$C_I+C_J$ contains all six multiples of $1/6$.  Given a target $t$, choose
$c_I\in C_I$, $c_J\in C_J$ and a signed displacement $\delta$ with
$|\delta|\leq 1/12$ such that
$t=c_I+c_J+\delta$ in $\mathbb T$.  Round
$c_I+\delta/2$ to a nearest grid point $i$, and put $j=t-i$.  Both $i$ and
$j$ are grid points, and their distances from $c_I$ and $c_J$, respectively,
are at most
\[
 \frac1{24}+\frac1{2q}
 <\frac1{24}+\frac1{800}<\frac1{16}.
\]
Moreover, the distance between any center in $C_I$ and any center in $C_J$
is at least $1/6$.  Consequently every allowed left-tag difference $i-j$
satisfies, in every batch coordinate,
\[
 \|i-j\|_{\mathbb T}\geq
 \frac16-\frac2{16}=\frac1{24}.
\]

For a right tag, any target batch with one small coordinate is a sum of
allowed batches.  To see this, write the batch size as $2k+1$ and choose a
small target coordinate.  At that coordinate, assign the target to the
first summand and zero to the second.  Partition the other $2k$ coordinates
into two sets of size $k$.  On the first set put zero in the first summand
and the target in the second; on the second set do the reverse.  Both
summands then have at least $k+1$ small coordinates.  Conversely, any two
allowed right-tag batches have a common small coordinate, because their
sets of small coordinates are strict majorities.  Their difference
therefore has some coordinate satisfying
\[
 \|i-j\|_{\mathbb T}\leq\frac2{100}<\frac1{24}.
\]
Thus the difference sets $I_\alpha-J_\alpha$ and
$I_{\alpha'}-J_{\alpha'}$ are disjoint whenever a block places $\alpha$
on the left and $\alpha'$ on the right.

\emph{The full relation, including independently mismatched tags.}
If $s\notin W$, all first coordinates of $s$ are nonzero and every batch
has a small coordinate.  The preceding sum constructions, together with
the first-coordinate decompositions $s=0+s$ and $s=s+0$, show that
$s\in I_\alpha+J_\alpha$ for every tag $\alpha$.  The choices in distinct
blocks are independent, so they combine to give elements of $G$.

Now suppose an off-diagonal pair $(\alpha,\alpha')$ belongs to
$\mathcal E(s)$.  Consider any competing pair $(\lambda,\mu)$ with
$\{\lambda,\mu\}\ne\{\alpha,\alpha'\}$.  There is an assignment
$\chi$ placing $\alpha$ on the left, $\alpha'$ on the right, and
$\lambda,\mu$ on the same side: identifying $\lambda$ with $\mu$ does
not identify $\alpha$ with $\alpha'$.  In its first coordinate the pair
$(\alpha,\alpha')$ forces $s=0$, whereas the pair $(\lambda,\mu)$
requires $s\ne0$.  Hence that competitor is impossible.  The only other
possible competitor is the reverse pair $(\alpha',\alpha)$.  If both
pairs represented $s$, there would be
\[
 i_\alpha+j_{\alpha'}=i_{\alpha'}+j_\alpha,
 \qquad\text{hence}\qquad
 i_\alpha-j_\alpha=i_{\alpha'}-j_{\alpha'},
\]
with the indicated elements belonging to their tag sets.  A block
separating these tags contradicts the disjointness of the difference
sets.  This proves the first alternative in \textup{(i)}.  It also excludes
off-diagonal pairs outside $W$, where the full diagonal is present, and
completes the proof of \textup{(i)}.

\emph{Assigning the roster and counting whole-batch obstructions.}
The tag relation is now established. What remains is a choice of the
$v_p$ that makes membership in $W$ rare among ordered roster differences,
for every common translate.  Write
$r=|\mathcal R|$ and enumerate the roster as $p_0,\ldots,p_{r-1}$.  In
every first coordinate set the residue of $v_{p_j}$ equal to $2^j$.  All
ordered differences for distinct roster slots are distinct modulo that
coordinate's prime.  For $r\geq2$, they are distinct as integers: the sign
determines their orientation, and in a positive difference
\[
 2^b-2^a=2^a(2^{b-a}-1),\qquad a<b,
\]
the power of two dividing it determines $a$, after which the odd factor
determines $b$.  Reduction modulo $q$ creates no collision, since the
distance between any two such integer differences is less than $2^r<q$.
The assertion is vacuous when $r<2$.  For any translate $u$, at most one
ordered pair therefore hits a specified first-coordinate zero obstruction.

In the dedicated batch coordinate for a list
$((p_a,p'_a))_{a=0}^{F_{\rm cl}-1}$, let $z_a$ be a grid point at circular
distance at most $1/(2q)$ from $a/F_{\rm cl}$.  Prescribe
\[
 (v_{p'_a}-v_{p_a})\big|_{\text{coordinate}}=z_a
 \qquad(0\leq a<F_{\rm cl}).
\]
These prescriptions can be satisfied: choose a root value in each tree
and assign successive vertex values along its edges, with the sign
determined by the prescribed orientation.  A forest has no consistency
condition around a cycle.  Assign arbitrary values to any remaining
vertices and in spare coordinates.  All these assignments are made
independently in each coordinate, and so define elements $v_p\in G$.

For any coordinate translate $t$, the equally spaced points
$a/F_{\rm cl}$ include one at distance at most $1/(2F_{\rm cl})$ from
$-t$.  The corresponding prescribed difference satisfies
\[
 \|t+z_a\|_{\mathbb T}
 \leq\frac1{2F_{\rm cl}}+\frac1{2q}
 <\frac1{400}<\frac1{100}.
\]
Thus every dedicated list has a pair whose translated difference is small
in its dedicated coordinate, uniformly in the translate.

Fix $u\in G$ and a batch.  Call an ordered pair $(p,p')$, $p\ne p'$,
bad for this batch if $u+v_{p'}-v_p$ has no small coordinate in the batch.
Suppose there were more than $4F_{\rm cl}^{2}$ such pairs.  Form the
simple undirected graph consisting of their underlying edges, discarding
isolated vertices.  If a spanning forest had fewer than $F_{\rm cl}$
edges, its nontrivial components would together have at most
$2F_{\rm cl}-2$ vertices: a component with $h\geq2$ vertices contributes
$h-1$ forest edges and satisfies $h\leq2(h-1)$.  The number of ordered
pairs would then be at most
\[
 (2F_{\rm cl}-2)(2F_{\rm cl}-3)<4F_{\rm cl}^{2},
\]
a contradiction.  Select $F_{\rm cl}$ edges of a spanning forest and
orient each as one of the available bad pairs.  In any order, they form a
list with a dedicated coordinate in this batch.  The preceding covering
estimate makes one of these pairs small in that coordinate, again a
contradiction.  Each batch therefore has at most $4F_{\rm cl}^{2}$ bad
ordered pairs.  This conclusion also applies when the roster is too small
to admit a dedicated list, by the same counting argument.

There are $2^{|\Lambda|}$ blocks.  Taking the union of their
first-coordinate and batch obstructions gives, for every $u\in G$,
\[
 \#\{(p,p'):p\ne p',\ u+v_{p'}-v_p\in W\}
 \leq 2^{|\Lambda|}\bigl(1+4F_{\rm cl}^{2}\bigr).
\]
For the empty roster the left side is zero throughout the construction.

\emph{One cyclic clock.}
Finally, let $n$ be the product of all coordinate primes.  The residue map
from $\mathbb Z/n\mathbb Z$ to $G$ is an injective homomorphism, since the
distinct primes are pairwise coprime, and both groups have $n$ elements.
It is therefore an isomorphism.  Transport all constructed sets and
residues through it.  This preserves sums, differences, and the two
properties, while every prime factor of $n$ exceeds $Q$.  All sets and
choices in the construction are finite.
\end{proof}

\section{Markers and several scales of cuts}\label{sec:scales}

The main split will use a bounded family of possible cuts.  The clock
controls which tags can meet at a cut, but a suffix still has several
possible origins for its computation.  We therefore choose boundaries
whose interval products remain unchanged under suitable movements of the
origin.  Markers keep a boundary fixed outside periodic regions; inside
such a region we may move it by a multiple of a sufficiently large
period.  Several scales of movement ensure a uniform bound on the
number of ambiguous origins.

There are three distinct outputs. The marker rule recognizes regions where
letters determine a common periodic extension. The roster and scale
construction provides many cut positions of every required kind. Finally,
product transport and a descending count show that only a bounded number
of ordered pairs of those positions can give obstructed origins. We develop
the finite data in that order and record exactly which later choices may
still be enlarged.

All parameters in this section are finite and may depend on $M$ and
$\Sigma$; lengths, positions, and thresholds are integers.  We first
work on two-sided words, indexed by the integers.
On a finite argument, every window prescribed by a rule must be present.
An unavailable window causes rejection; in particular, it never removes
an origin from a universal or uniqueness test.  Section~\ref{sec:episodes}
will place all the required windows inside complete episodes.

\subsection{Local markers and periodic extensions}

A positive integer $a$ is a \emph{period} of a finite word if its letters
at distance $a$ agree whenever both positions belong to the word.

The marker-versus-local-periodicity viewpoint goes back to Krieger's marker
lemma~\cite[Lemma~2]{Krieger1982}; see also the finite clopen merging
formulation in Meyerovitch~\cite[Section~3, Lemmas~3.2--3.4]{Meyerovitch2025}.
We prove the finite-window rule needed here. The long-block hypothesis
$K>3d^2$ also prepares the overlap argument in the next lemma.

\begin{lemma}[Local markers]\label{lem:markers}
Let $d\geq 2$ and $K>3d^2$ be integers.  There is a
translation-equivariant rule marking integer positions of any two-sided
word such that distinct markers are at distance at least $d$.
The decision at $z$ uses only $[z-r,z+r)$, for a fixed finite integer
$r\geq K$.  If there is no marker in the inclusive integer interval
$[z-d,z+d]$, then the block $[z,z+K)$ has a positive period less than $d$.
\end{lemma}

\begin{proof}
List the finitely many length-$K$ words having no positive period below
$d$.  Process these words in a fixed order.  When processing a word,
mark all its starts that are not within distance less than $d$ of a
previously installed marker.  Two starts added at this same step cannot
be at distance $a<d$: their overlap would make $a$ a period of the
listed word.  Induction therefore proves separation of all markers.

Each step makes its decision from a length-$K$ block and the decisions
of earlier steps at offsets of magnitude less than $d$.  Induction over
the finite list gives a finite inspection radius; enlarge it to at
least $K$.  The construction commutes with translations.  Finally, a
listed word at $z$ was either marked or prevented by a marker within
distance less than $d$.  Thus marker absence on $[z-d,z+d]$ excludes
every listed word at $z$, as required.
\end{proof}

\begin{lemma}[Common periodic extensions]\label{lem:periodic-extension}
A block of length $K>3d^2$ having a positive period below $d$ determines
a unique two-sided periodic extension of period below $d$.  Two such
blocks whose starts differ by less than $K-d^2$ determine the same
extension.  For the marking rule of Lemma~\ref{lem:markers}, if there is
no marker in $[z-H,z+H]$ and $H\geq d$, the actual letters on
\[
       [z-H+d,\ z+H-d+K)
\]
belong to one such periodic extension.
\end{lemma}

\begin{proof}
A block of period $a$ extends by repeating its first $a$ letters.  Two
two-sided words of respective periods $a,b$ that agree on $ab$
consecutive positions agree everywhere: both have period $ab$.
Here $a,b<d$, so $ab<d^2$.  This proves both uniqueness and the assertion
about overlapping blocks.

For the last assertion, every start in the inclusive interval
$[z-H+d,z+H-d]$ has no marker within distance $d$.  Its length-$K$ block
therefore has a period below $d$.  Consecutive such blocks overlap
sufficiently, so their extensions agree.  Their union is the displayed
interval, on which the extension agrees with the actual letters.
\end{proof}

Whenever these lemmas are applied to a finite word, extend the word
arbitrarily in both directions.  If all prescribed inspection windows
are internal, the computed markers and periodic blocks are independent
of that extension.

\subsection{Algebraic data and the order of the buffers}

Before choosing cut positions, we fix the states and finite information
that each cut must transmit. Their cardinalities determine how many
candidate cut positions we need.

Put
\[
 U_0=V,\qquad U_{i+1}=\F_2^{\,U_i\times M}\quad (i=0,1).
\]
We denote the basis vector of $U_{i+1}$ indexed by $(v,c)$ by $[v,c]$.
Its work action is
\begin{equation}\label{eq:work-action}
          [v,c]\rho_i(d)=[v,cd]\qquad(d\in M).
\end{equation}
Thus the formal label $[v,c]$ retains an input state $v$ while its
second coordinate accumulates a prefix product. Once the later part of
an episode is known, a terminal linear map will send each label to a
basis label recording the completed update corresponding to $c$.
Lemma~\ref{lem:table-lift} constructs that map for every $c\in M$,
including values not realized by an actual prefix. The two larger spaces
will let us perform this encoding twice.

Choose $l,\beta$ from Lemma~\ref{lem:affine-line} for the action of $M$
on $V$.  For each $i=0,1$, choose $l_i,\beta_i$ from the same lemma for
the action $\rho_i$ on $U_{i+1}$.  These choices are made once and for all.
Thus $\beta:M^l\to V$ and $\beta_i:M^{l_i}\to U_{i+1}$.

An episode beginning at $s$ will have a later origin
\[
                         t=s+L.
\]
Before $t$ we will place two buffers in the order $1,0$.  Buffer $i$
has checkpoints
\[
 B_{i,0}<B_{i,1}<\cdots<B_{i,l_i}<a_i,
\]
where $a_i$ is after all its activities.  The entire buffer $1$,
including $a_1$, precedes buffer $0$.  The checkpoints are translates
of fixed offsets from $s$; we use $B_{i,j}$ and $a_i$ for their absolute
positions in the word.  The translations and $L$ will be fixed in
Section~\ref{sec:episodes}.

The \emph{kind} of a cut in buffer $i$ is a tuple
\begin{equation}\label{eq:buffer-kinds}
 (j,v,d_1,\ldots,d_{j-1}),\qquad
 1\leq j\leq l_i,\quad v\in U_{i+1},\quad d_f\in M.
\end{equation}
Let $\kappa_i$ be the number of these kinds.  A finite stencil
$\mathcal P_i$ will specify the cut offsets $p$ from $s$, with each
offset assigned a kind.  An offset of interval number $j$ must satisfy
$B_{i,j-1}<s+p<B_{i,j}$.  Write
\begin{equation}\label{eq:buffer-differences}
 h_i=\operatorname{diam}\mathcal P_i,\qquad
 \Delta_i=\mathcal P_i-\mathcal P_i.
\end{equation}
At an actual cut $s+p$, a suffix considering an offset $p'$ will infer
the origin $t+p-p'$.  We construct $\mathcal P_0$ here and
$\mathcal P_1$ in Section~\ref{sec:episodes}.

Once separated numerical offsets have been chosen, kinds can be
assigned in increasing interval-number order.  Boundaries between
successive interval numbers are then placed in the gaps between their
offsets; the extreme boundaries and $a_i$ can be added outside them.
Thus it suffices to produce separated offsets with sufficiently many
copies of each kind.

\subsection{The clock tags and a polynomial roster}

The first stencil must contain more copies of every kind than the total
number of wrong ordered origin pairs that can survive our tests. The
following polynomial roster leaves this growth margin: its number of slots
grows by a higher power of the field size than the nonclock obstruction
budget does. The exponents and other constants are fixed before that
field size is selected.

For the main segment set
\begin{equation}\label{eq:scale-counts}
             Z=l+1,\qquad k_0=2Z+1,\qquad R=2k_0-1.
\end{equation}
There will be $Z$ centers for its boundaries and $R$ scales for each
cut role.  Use the following tag set in Lemma~\ref{lem:clock}:
\begin{equation}\label{eq:clock-tags}
 \alpha=(j,r,v,d_1,\ldots,d_{j-1},X,Y),\qquad
 \begin{gathered}
 1\leq j\leq l,\quad 1\leq r\leq R,\\[-2pt]
 v,X,Y\in V,\quad d_f\in M.
 \end{gathered}
\end{equation}
Its ambiguity constant $C_{\rm cl}$ is determined by the tag set alone,
before a roster or its positions have been chosen.

Choose a prime $P>R$ large enough that
\begin{equation}\label{eq:stage0-constant}
 \left\lfloor\frac{P^{k_0}}{\kappa_0}\right\rfloor>C_0,
 \qquad
 C_0=C_{\rm cl}+2^R(3Z+1)^R P^{2Z}.
\end{equation}
Such primes exist because $k_0=2Z+1>2Z$.  The abstract roster consists
of all polynomials of degree less than $k_0$ over $\F_P$, including the
zero polynomial. This is the polynomial-evaluation construction underlying
Reed--Solomon codes~\cite{ReedSolomon1960}; we use only the elementary
root bound, which applies over the prime field chosen here.
Choose distinct $\xi_1,\ldots,\xi_R\in\F_P$ and give a
roster slot $p$ the digits
\[
                         x_r(p)=p(\xi_r).
\]
We identify field elements with $\{0,\ldots,P-1\}$ when using them as
digit indices.  A nonzero polynomial of degree less than $k_0$ has at
most $k_0-1$ roots, by successive division by linear factors.
Consequently any $k_0$ evaluations determine a slot, and two distinct
slots differ in at least
\begin{equation}\label{eq:roster-distance}
                         R-k_0+1=k_0
\end{equation}
digits.

Apply Lemma~\ref{lem:clock} to this roster, avoiding all primes at most
$N$.  Fix its modulus $n$, tag sets $I_\alpha,J_\alpha$, bad set $W$,
and roster residues $v_p\in\mathbb Z/n\mathbb Z$.  The prime divisors
of $n$ are all greater than $N$.  Notice that the modulus is chosen
after the roster size; enlarging it does not change $C_0$.

\subsection{Scale choices and the first stencil}

The abstract roster is now fixed, as are its required clock residues. We
next realize the slots at integer positions. Large separation at each
scale will dominate all smaller-digit changes, while a small residue
correction preserves the clock assignment.

Call a positive integer \emph{smooth} if all its prime divisors divide
$N$.  We choose thresholds $S_r<S'_r$, a positive step $\delta_r$,
a movement bound $R_r=n\delta_r$, and digit positions
$f_r(0)<\cdots<f_r(P-1)$, in increasing order of $r$.
For each $r$ perform the following choices:
\begin{enumerate}[label=(\roman*)]
\item Choose $S_r\geq1$, also requiring $S_r\geq S'_{r-1}$ if $r>1$.
Choose $\delta_r$ smooth and divisible by $N\sigma$ for every smooth
positive integer $\sigma\leq S_r$.  Then
\begin{equation}\label{eq:step-coprime}
                            \gcd(\delta_r,n)=1.
\end{equation}
\item If $h'_b$ denotes the diameter of the digit positions at scale
$b$, put
\[
                    J_r=n+\sum_{b<r}h'_b.
\]
Choose the digit positions so that all ordered differences of unequal
digits are distinct and separated from one another and from zero by
more than $10(J_r+R_r+1)$.  Denote their diameter by $h'_r$.
\item Choose
\begin{equation}\label{eq:upper-scale-threshold}
              S'_r>\max\{S_r,\,10(h'_r+J_r+R_r+1)\}.
\end{equation}
\end{enumerate}
The step in (i) may be the least common multiple of the finitely many
required $N\sigma$; it is smooth.  If $N=1$, the only smooth positive
integer is $1$, and $\delta_r=1$ is permissible.  The digit positions
in (ii) can be sufficiently scaled powers of two.  Indeed a positive
difference $2^a-2^b$, $a>b$, determines $b$ by its power-of-two
valuation and then determines $a$.  All nonzero ordered differences
are thus distinct, and scaling provides the required separation.

Assign slot $p$ the numerical offset
\begin{equation}\label{eq:slot-position}
                   \sum_{r=1}^R f_r(x_r(p))+\epsilon_p,
             \qquad 0\leq\epsilon_p<n.
\end{equation}
Choose $\epsilon_p$ to give residue $v_p$ modulo $n$.
For two slots whose highest differing digit is at scale $r$, the sum
of all smaller-coordinate differences and the correction difference
has absolute value less than $J_r$.  The difference at scale $r$
has absolute value greater than $10(J_r+R_r+1)$.  The resulting
positions are therefore distinct, and every two are separated by
more than $2$.

These positions, up to a common translation, form $\mathcal P_0$.
We henceforth identify a slot with its numerical offset when taking
differences.  Their difference residues are exactly the prescribed
$v_p-v_{p'}$, including after the common translation.  By
\eqref{eq:stage0-constant}, we can assign more than $C_0$ offsets to
each stage-$0$ kind.  Assign them in interval-number order and insert
the buffer boundaries as described above.  Translation of the complete
configuration later changes none of its difference properties.  In
particular,
\begin{equation}\label{eq:stage0-diameter}
                         h_0<n+\sum_r h'_r.
\end{equation}

\subsection{Inner controls and the main boundaries}

We have fixed the stencil and every movement range. We can now make the
inner markers sparse relative to all of them. In marker-free regions the
boundary rule must preserve a congruence class modulo a full monoid period;
choosing only a phase modulo the word period would not suffice.

Choose integers
\begin{equation}\label{eq:inner-marker-size}
 d_\circ>10\left(1+n+\sum_r h'_r+\sum_r R_r\right),
 \qquad K_\circ>3d_\circ^2,
\end{equation}
and fix the marking rule of Lemma~\ref{lem:markers} with these
parameters and radius $r_\circ\geq K_\circ$.  Choose an integer $H$ with
\begin{equation}\label{eq:inner-halfwidth}
             H>10\bigl((N+m+1)d_\circ+\max_r R_r+1\bigr).
\end{equation}
For an inner center $z$, let $u(z)$ be the first marker in
$[z-H,z+H]$, or $\bot$ if none exists.  If $u(z)\neq\bot$, put
$B(z)=u(z)$.

If $u(z)=\bot$, Lemma~\ref{lem:periodic-extension} supplies a periodic
extension of the block starting at $z$.  Let $d<d_\circ$ be its least
period, and let $\sigma$ be the largest smooth divisor of $d$.  Among
the length-$d$ period words starting at its different phases, choose
the lexicographically least, using a fixed order on $\Sigma$.  Its
starts form one class $x\bmod d$: equality of two rotations would
contradict minimality of $d$.  Define $y(z)$ to be the largest integer
at most $z$ in $x+\sigma\mathbb Z$, and let $B(z)$ be the least integer
at least $z$ satisfying
\begin{equation}\label{eq:boundary-congruences}
                  B(z)\equiv x\pmod d,\qquad
                  B(z)\equiv y(z)\pmod{\sigma N}.
\end{equation}
The definition is independent of the representative of $x$.  Moreover,
\[
 \gcd(d,\sigma N)=\sigma,\qquad
 \operatorname{lcm}(d,\sigma N)=dN,
\]
because $d/\sigma$ is coprime to $N$.  The two residues in
\eqref{eq:boundary-congruences} agree modulo $\sigma$, so they define
one class modulo $dN$.  In particular,
\begin{equation}\label{eq:periodic-boundary-range}
                            z\leq B(z)<z+dN.
\end{equation}
All these choices commute with translations of the letters.  They
do not use an external choice of position origin.  In both the present
and absent cases, $B(z)\in[z-H,z+H]$.

Choose an even integer
\begin{equation}\label{eq:center-spacing}
                    G>10(H+r_\circ+\max_r R_r+1).
\end{equation}
For the main segment set
\[
 o_j=(j+1)G,\quad z_j=t+o_j,\quad B_j=B(z_j)
       \quad(0\leq j\leq l),
 \qquad c_j=\frac{o_{j-1}+o_j}{2}\quad(1\leq j\leq l).
\]
At another proposed origin $\tau$, write
$z_j(\tau)=\tau+o_j$ and $B_j(\tau)=B(z_j(\tau))$; the notation
without an argument refers to the true origin $t$.
The main cut offsets from $t$ form the finite set
\begin{equation}\label{eq:main-cuts}
 \mathcal H=
 \{c_j+\delta_r a:1\leq j\leq l,\ 1\leq r\leq R,\ 0\leq a<n\}.
\end{equation}
A cut of role $(j,r)$ lies strictly between $B_{j-1}$ and $B_j$.
All inspection windows for boundaries through $j-1$ are before this
cut.  Conversely, for a suffix beginning at $k$ and any hypothesized
role $(j,r)$, a candidate origin $k-c_j-\delta_r a$ has its first
required center at
\begin{equation}\label{eq:suffix-first-center}
                       k+G/2-\delta_r a.
\end{equation}
Since $\delta_r a<R_r$, \eqref{eq:center-spacing} puts the inspections
for this center and all later centers after $k$.  This holds even when
the suffix role differs from the prefix role.  Complete placement of
these controls, including the later end controls, is deferred to
Lemma~\ref{lem:placement}.

\subsection{Usable scales and transport of products}

The boundary rule is defined at every origin. We next specify when a
whole range of shifted origins gives equal interval products. The
condition compares actual marker positions, including absence, and in
periodic regions controls the phase lattice used by the congruence rule.

\begin{definition}\label{def:usable}
A scale $r$ is \emph{usable at the origin $t$} if, at each of its $Z$
centers $z_j$, the following conditions hold:
\begin{enumerate}[label=(\roman*)]
\item The value $u(z_j+h)$ equals $u(z_j)$ for every integer $h$ with
$|h|\leq R_r$, including agreement on absence.
\item If $u(z_j)=\bot$, then either $\sigma\leq S_r$, or the inclusive
interval $[z_j-R_r,z_j+R_r]$ contains no point of
$x+\sigma\mathbb Z$.
\end{enumerate}
Here $d,\sigma,x$ refer to the periodic extension at the indicated
center.  All windows needed by these checks are mandatory.
\end{definition}

\begin{lemma}[Transport of boundary products]\label{lem:boundary-transport}
Suppose $r$ is usable at $t$, and $h$ is a multiple of $\delta_r$ with
$|h|\leq R_r$.  Write $B_j^{(h)}=B(z_j+h)$.  Then
\begin{equation}\label{eq:transport-intervals}
        T_{B_{j-1},B_j}=T_{B_{j-1}^{(h)},B_j^{(h)}}
                         \qquad(1\leq j\leq l).
\end{equation}
For a fixed endpoint $b\geq z_l+H$, the last transfer also satisfies
\begin{equation}\label{eq:transport-tail}
                         T_{B_l,b}=T_{B_l^{(h)},b}.
\end{equation}
These statements apply to finite words whenever the prescribed
windows and product intervals are available.  They hold for arbitrary
finite monoids, without a commutativity assumption.
\end{lemma}

\begin{proof}
The proof has three steps: the two boundaries stay in one class modulo
$dN$; the actual word contains enough periodic letters on both sides;
and monoid-power periodicity then compares the ordered products.

\emph{A common boundary class.}
Consider one center $z$. If its marker is present, usability leaves
the boundary unchanged.  If absent, the old and new periodic blocks
have the same extension by Lemma~\ref{lem:periodic-extension}:
$|h|\leq R_r<d_\circ$, and the blocks have length $K_\circ>3d_\circ^2$.
Thus the least period $d$, smooth divisor $\sigma$, and phase class
$x\bmod d$ are common.

If $\sigma\leq S_r$, the displacement $h$ is divisible by $\sigma N$.
Translation along the lattice gives $y(z+h)=y(z)+h$, so the second
residue in \eqref{eq:boundary-congruences} is unchanged.  Otherwise,
usability says that no point of that lattice lies between any of the
positions under consideration, so $y(z+h)=y(z)$.  In both cases the
two boundaries belong to one class modulo $dN$.  Writing them in
increasing order as $b_-\leq b_+$, we have
\begin{equation}\label{eq:boundary-displacement}
                            b_+-b_-=kdN\quad(k\geq0).
\end{equation}

\emph{Margins of actual periodic letters.}
By \eqref{eq:periodic-boundary-range}, both choices lie in
$[z-R_r,z+R_r+Nd)$.  Marker absence and
Lemma~\ref{lem:periodic-extension} give a common run of actual letters
containing $[z-H+d_\circ,z+H-d_\circ)$.  The inequality
\eqref{eq:inner-halfwidth} therefore implies that
\begin{equation}\label{eq:periodic-margins}
               [b_--md,\ b_++md)
                 \subset [z-H+d_\circ,z+H-d_\circ)
\end{equation}
lies in that run.  In particular both choices have the required
periodic margins on their two sides.

\emph{Comparing products one endpoint at a time.}
Let $c\in M$ be the product of one length-$d$ period beginning in
the phase of $b_-$.  For an incoming product from any fixed
$a\leq z-H$, comparison at $b_--md$ gives the two factorizations
\[
             T_{a,b_-}=A c^m,\qquad
             T_{a,b_+}=A c^{m+kN}.
\]
For an outgoing product to any fixed $b\geq z+H$, comparison at
$b_++md$ gives
\[
             T_{b_-,b}=c^{m+kN}C,\qquad
             T_{b_+,b}=c^m C.
\]
Here $A$ and $C$ are the products on the unchanged remaining portions.
Equation~\eqref{eq:monoid-periodicity} gives
$c^{m+kN}=c^m$.  The equalities require neither cancellation nor
reordering of factors.

The old and new boundary at each center lie within its unshifted
window $[z_j-H,z_j+H]$, by the same estimates.  Adjacent center
windows are disjoint and separated by more than their required
margins, since $G>10H$.  To compare an interval between successive
boundaries, move its left endpoint while fixing its right endpoint,
then move the right endpoint while fixing the new left endpoint.
Each fixed opposite endpoint is outside the relevant window, so the
local calculation applies.  This proves
\eqref{eq:transport-intervals}.  Its outgoing version at the last
center proves \eqref{eq:transport-tail}.
\end{proof}

\subsection{A uniform ambiguity bound for the first stencil}

Product transport handles an origin whenever at least one scale is usable.
We now bound all the pairs of stencil slots that can produce an origin
where this fails, also allowing a bad clock residue. Throughout this count
the anchor $q$ is fixed. Permitting it to vary with a candidate would not
give the translated-difference problem solved by the clock.

For a proposed origin $t$ and an anchor position $q$, define
\begin{equation}\label{eq:stage0-predicate}
 A_0(t,q):\qquad
     q-t\pmod n\in W
       \quad\text{or no scale is usable at }t.
\end{equation}
The anchor is held fixed when comparing different proposed origins.

\begin{proposition}[Inner ambiguity]\label{prop:inner-ambiguity}
On any two-sided word and for any integers $t,q$,
\begin{equation}\label{eq:inner-ambiguity}
 \#\{(p,p')\in\mathcal P_0^2:p\neq p',\
                         A_0(t+p-p',q)\}\leq C_0.
\end{equation}
The same bound holds on finite words whenever all the indicated
inspection windows are internal.
\end{proposition}

\begin{proof}
We count numerical difference tuples rather than choose digit pairs at
every stage. Nonzero differences recover the slots; zero differences
contribute only one value. This distinction prevents an unnecessary factor
of $P$ at each equal-digit coordinate.

\emph{Clock failures and recovery of the ordered pair.}
The clock part contributes at most $C_{\rm cl}$, because
\[
      q-(t+p-p')\equiv(q-t)+v_{p'}-v_p\pmod n.
\]
This is precisely the uniform translated-difference bound of
Lemma~\ref{lem:clock}.  It remains to count pairs for which no scale
is usable.

For an ordered pair of slots define its numerical digit differences
\[
                  D_r=f_r(x_r(p))-f_r(x_r(p')).
\]
The full tuple $(D_1,\ldots,D_R)$ determines a distinct ordered pair
$p\neq p'$.  Indeed, it has at least $k_0$ nonzero coordinates by
\eqref{eq:roster-distance}.  Each such coordinate determines both
ordered digits by the uniqueness of the digit differences.  These
$k_0$ evaluations determine each polynomial slot separately; unrevealed
equal digits at zero coordinates therefore introduce no additional pairs.

\emph{Exceptional scales.}
For a pair with no usable scale, look at its origin $t+p-p'$ and
declare a coordinate $r$ exceptional if
\begin{equation}\label{eq:exceptional-scale}
                S_r<\sigma\leq S'_r
\end{equation}
at at least one absent center.  At this fixed origin each center has
one value of $\sigma$, and the intervals $(S_r,S'_r]$ are disjoint.
There are therefore at most $Z$ exceptional coordinates.  We may sum
over all their possible sets, of which there are at most $2^R$.

\emph{Descending count with all smaller digits still free.}
Fix such a set of exceptions.  Count difference tuples in descending
coordinate order.  At an exceptional coordinate there are at most
$P^2$ possible differences.  At a nonexceptional coordinate $r$, fix
the differences at all larger coordinates.  For each center label
$j$, its possible positions are of the form
\begin{equation}\label{eq:center-digit-error}
 C_j+D_r+e,\qquad |e|\leq J_r,\qquad
 C_j=t+o_j+\sum_{b>r}D_b.
\end{equation}
The error includes the smaller-coordinate differences and
$\epsilon_p-\epsilon_{p'}$.  This description fixes only the larger
numerical differences; it does not assume that their underlying
digit pairs have already been determined.

Since scale $r$ fails to be usable, at some center label one of the
following occurs:
\begin{enumerate}[label=(\alph*)]
\item A marker lies within $R_r$ of one of the two search endpoints.
\item The search is absent at the center, $\sigma>S'_r$, and a point
of $x+\sigma\mathbb Z$ lies within $R_r$ of the center.
\end{enumerate}
Indeed a change in the first marker of a translated search interval
requires a marker to enter or leave through an endpoint.  If the
search does not change, failure of the second usability condition
requires absence, $\sigma>S_r$, and such a lattice point.
Nonexceptionality then strengthens $\sigma>S_r$ to $\sigma>S'_r$.

For (a), the range swept by a given search endpoint, enlarged by
$R_r$, has diameter at most
\begin{equation}\label{eq:swept-range}
                  2(h'_r+J_r+R_r)<d_\circ.
\end{equation}
It contains at most one marker.  A fixed marker can serve at most
one nonzero value of $D_r$: by \eqref{eq:center-digit-error}, serving
it restricts $D_r$ to an interval of length $2(J_r+R_r)$, whereas
distinct nonzero digit differences are more widely separated.
The two endpoints give at most two such values for each label.

For (b), all the centers under consideration for a fixed label lie
in an interval of diameter at most $2(h'_r+J_r)<d_\circ$.
Among those centers whose searches are absent, their length-$K_\circ$
blocks determine the same two-sided periodic extension, by
Lemma~\ref{lem:periodic-extension}: their start distance is less than
$d_\circ$, and $K_\circ-d_\circ>d_\circ^2$. Only absent candidates enter
this comparison. Present candidates between them are irrelevant, since
overlap directly compares any two of the absent candidates' seed blocks.
Thus, even as $D_r$, the lower digits, and the residue corrections vary, these absent candidates have one
common least period, smooth divisor, phase, and lattice.  This is
where the long local blocks are needed in the descending count.
If the common $\sigma$ does not exceed $S'_r$, none of these candidates
can realize (b).  Otherwise the enlarged center range has diameter
less than $\sigma$, by \eqref{eq:upper-scale-threshold}; it contains
at most one point of the common lattice.  That point accounts for
at most one nonzero difference $D_r$, by the same separation argument
as for a marker.

\emph{Summing the budgets.}
There are $Z$ center labels.  Allowing also the single difference
value zero, a nonexceptional coordinate therefore has at most
$3Z+1$ possibilities.  It is not necessary to specify the equal
digits when the difference is zero.  Summing the descending count
over exception sets bounds the number of full difference tuples by
\[
                  2^R P^{2Z}(3Z+1)^R.
\]
Since different ordered pairs give different tuples, this also bounds
the number of pairs.  Adding the clock contribution proves
\eqref{eq:inner-ambiguity} with the constant in
\eqref{eq:stage0-constant}.

For a finite word, take any two-sided extension.  Internal inspection
windows give the same decisions, so the bound just proved applies
unchanged.
\end{proof}

The order of all choices is worth recording.  The affine data and tag
set fix $C_{\rm cl}$; then $P$ fixes the roster; then the clock fixes
$n$; then the scale parameters and digit positions are chosen
successively.  Only after these choices do we choose the inner marker
parameters, $H$, and $G$.  Thus neither the ambiguity estimate nor the
number of copies per kind depends circularly on a later geometric
choice.  The remaining buffer, end controls, and translations can now
be added around these fixed finite arrangements.

\section{Episodes and exact end detection}
\label{sec:episodes}

The first ambiguity estimate held its anchor fixed. We must now arrange
that independently accepted suffix tests actually use that same anchor.
This section builds an episode rule and proves the necessary exact-end
statement. The main split and the first residual split compare a full list
of bounded marker searches, then run one unbounded end search at a reference
origin. The second residual split uses a different order: bounded tests
decode an origin before its one end search. Keeping these orders distinct
is necessary both for soundness and for the later height-one compilation.

We first construct the second stencil and bound instability of its end
search. We then place every window, including windows at wrong candidate
origins, inside a complete episode. The final guard lemma treats arbitrary
accepted suffixes, not merely the intended suffix of that episode.

\subsection{The second stencil and the end obstruction}

Retain the stage-0 stencil and main cut set from Section~\ref{sec:scales}.
Choose an integer
\[
 w>\max\{h_0,\operatorname{diam}\mathcal H\},
 \qquad C_1=2|\Delta_0|.
\]
For each stage-1 kind, take more than $C_1$ copies.  Give the copies
distinct integer offsets whose nonzero ordered differences are all distinct
and are separated from one another and from zero by more than $2w+10$.
Scaled powers of two provide such offsets: the positive difference
$2^b-2^a$, $b>a$, determines $a$ from its exact power of two and then
determines $b$; a sufficiently large common scale gives the separation.
Assign the kinds to the offsets in interval-number order and insert the
buffer boundaries in the intervening gaps, as in
Section~\ref{sec:scales}.  Denote the resulting stencil by $\mathcal P_1$,
and put
\[
 h_1=\operatorname{diam}\mathcal P_1,
 \qquad \Delta_1=\mathcal P_1-\mathcal P_1.
\]
We will translate each entire buffer arrangement later; translations do not
affect these difference properties.

Choose integers $d_e,K_e$ with
\[
 d_e>10(1+w+h_1),\qquad K_e>3d_e^2,
\]
and apply Lemma~\ref{lem:markers} to obtain an end-marker rule of radius
$r_e\ge K_e$.  An end-search center will have the form
$z_e(t)=t+o_e$, where $o_e$ is chosen below.  Let $\eta(t)$ be the first
end marker in the inclusive interval
$[z_e(t)-d_e,z_e(t)+d_e]$, or $\bot$ if none exists.  Define the bounded
predicate
\begin{equation}
 \label{eq:end-obstruction}
 A_1(t)\quad\Longleftrightarrow\quad
 \eta(t+u)\ne\eta(t)\text{ for some integer }|u|\le w.
\end{equation}
Present markers are compared as positions in the word, not as offsets from
their respective search centers.

\begin{lemma}[End ambiguity]
\label{lem:end-ambiguity}
On a two-sided word, uniformly in $t$,
\begin{equation}
 \label{eq:end-ambiguity}
 \#\bigl\{(p,p')\in\mathcal P_1^2:p\ne p',\
       \exists\gamma\in\Delta_0\ A_1(t+p-p'+\gamma)\bigr\}
 \le C_1.
\end{equation}
The same estimate holds on a finite word whenever all the indicated
inspection windows are present.
\end{lemma}

\begin{proof}
If the first marker in a moving search interval changes under a shift of
magnitude at most $w$, a marker must cross one of its endpoints.  In
particular, a marker lies within distance $w$ of an original endpoint.
Fix $\gamma\in\Delta_0$ and fix one of the two endpoints.  As $p-p'$
ranges over $\Delta_1$, the possible marker positions just described lie
in an interval of diameter at most $2h_1+2w<d_e$.  This interval contains
at most one end marker.  For that marker, the required ordered difference
$p-p'$ lies in an interval of length $2w$.  The separation of the nonzero
ordered differences gives at most one ordered pair with $p\ne p'$.
There are two endpoints and $|\Delta_0|$ choices of $\gamma$, proving
the bound.  With all windows available in a finite word, take any
two-sided extension; none of the inspected marker decisions changes.
\end{proof}

\subsection{Placement and the episode rule}

The end rule must allow arbitrarily long episodes while exposing just one
unbounded continuation. It does so by following a short periodic seed to
its first mismatch. A fixed tail after the anchor leaves room for every
bounded control window, even if that mismatch occurs immediately.

At a start $s$ with origin $t=s+L$, the prospective anchor is defined as
follows.  If $\eta(t)$ is present, set $q=\eta(t)$.  Otherwise the seed
\[
 [z_e(t),z_e(t)+K_e)
\]
has a period less than $d_e$ by Lemma~\ref{lem:markers}, and a unique
two-sided periodic extension by Lemma~\ref{lem:periodic-extension}.
Set $q$ equal to the first letter position after the seed at which the word
differs from this extension, if such a position exists.  For a fixed
positive integer $b_{\rm tail}$, the prospective episode end is
$q+b_{\rm tail}$.  An absent marker and no mismatch yet do not define a
complete episode.

We use the following convention in every finite-argument test in the
paper.  A letter inspection requires its entire window to lie in the
argument; a product requires its two endpoints to be available and
ordered.  A universal or uniqueness test on a finite candidate list
requires the bounded control windows for \emph{every} candidate.  If any
such window is missing, the test rejects; its candidate list is never
shortened.  A later computation at a decoded candidate also requires its
own data.  These requirements are finite checks, and are part of the
languages being defined.

\begin{lemma}[Placement]
\label{lem:placement}
The translations of the two buffers and the integers $L,o_e,b_{\rm tail}$
can be chosen so that the following properties hold.
\begin{enumerate}[label=\textup{(\roman*)}]
\item After $s$, the buffers occur in the order $1,0$, followed by all
relevant inner-system windows and cuts, and then by all relevant end-system
windows and seeds.
\item All bounded windows required at main hypotheses
$k-h$, $h\in\mathcal H$, or at aligned hypotheses
$k+L-p'$, $p'\in\mathcal P_i$, are internal to every complete episode
at a true cut.  This includes all usability checks, the additional
$\gamma\in\Delta_0$ translates at stage~1, and all shifts in $A_1$.
\item A prefix or suffix computation has its required data on its own
side of a true cut.  Later main data for any guessed role, and later
buffer data for any decoded offset, are scheduled after the suffix start.
\item Every possible true first-episode end is after all these cuts and
bounded control windows.
\end{enumerate}
\end{lemma}

\begin{proof}
\emph{Three finite envelopes.}
First translate the fixed buffer arrangements intact so that
\[
 s<B_{1,0}<a_1<B_{0,0}<a_0
\]
with each buffer's own boundaries and cut offsets in their prescribed
order.  Each difference $B_{i,j}-s$ or $a_i-s$ is a fixed offset.
Set
\[
 D=\operatorname{diam}\mathcal H+h_0+h_1+w+1.
\]
The main hypotheses differ from the true origin by at most
$\operatorname{diam}\mathcal H$.  Aligned hypotheses differ by
$p-p'$, of magnitude at most $h_i$; stage~1 adds $\gamma$ of magnitude
at most $h_0$, and an $A_1$ test adds at most $w$.  Thus $D$ bounds all
the end-system displacements that will be used.

Include in an inner envelope every inner search, inspection block,
boundary, and main cut at origins displaced by at most $D$, together with
the additional shifts of magnitude at most $\max_r R_r$ needed for
usability.  Every such position lies in $[t-C,t+C]$ for a finite constant
$C\ge D+1$ fixed before $L$ is chosen.  The marker radii and the bounded phase
displacements used to define the $B_j$ are included in $C$.  Choose $L$
so that $L-C>a_0-s$.

Likewise, include all end searches, their marker-inspection windows, and
their seeds at displacements of magnitude at most $D$.  They lie in
$[z_e(t)-C',z_e(t)+C']$ for a finite constant $C'$ independent of $o_e$;
increase it so that $C'>d_e+D$.  Choose $o_e$ so that $o_e-C'>C$.
Finally choose
\[
 b_{\rm tail}>C'+d_e+1.
\]
Any anchor satisfies $q\ge z_e(t)-d_e$.  Therefore
$q+b_{\rm tail}>z_e(t)+C'+1$, proving that the true end is beyond the
whole end envelope, and hence beyond both earlier envelopes.

\emph{Data on the required side of each cut.}
It remains to verify the sidedness of the computations.  At a main cut
$k=t+c_j+\delta_r a$, the spacing of the main centers puts all
inspections for boundaries through $B_{j-1}$ before $k$.  For any
guessed role $(j,r)$ in a suffix and any corresponding hypothesis
$t'=k-c_j-\delta_r a$, its first required center is
\[
 z_j(t')=k+G/2-\delta_r a.
\]
Since $0\le\delta_r a<R_r$ and
$G>10(H+r_\circ+\max_r R_r+1)$, its inspection windows, boundary, and
all later main data lie after $k$.  This calculation holds independently
of the role that produced the actual prefix cut.

For a buffer hypothesis $p'$ of interval kind $j$, the hypothetical
start is $s'=k-p'$.  The first later boundary has offset
$(B_{i,j}-s)-p'>0$ from $k$, and all required preceding data lie on the
prefix side when the hypothesis is true.  The order $1,0$ puts all
stage-0 segment data after every stage-1 cut; hence the stage-1 suffix can
form its later terminal table.  All inner and end data are later still
by the envelope choices.  This proves every assertion.
\end{proof}

All the choices in this proof are acyclic.  The inner construction fixes
$\mathcal P_0$ and $\mathcal H$ before $w$ is chosen; then
$\mathcal P_1,d_e,K_e,r_e$ are fixed.  Translating the buffers leaves
their differences unchanged.  Only after that do we choose $L$, then
$o_e$, then $b_{\rm tail}$.  None of these last choices changes an
ambiguity bound.

\begin{figure}[tbp]
 \centering
 \begin{tikzpicture}[x=0.91cm,y=0.75cm,font=\small]
 \draw[->,thick] (0,0) -- (15,0);
 \draw (0,0.12) -- (0,-0.12) node[below=3pt] {$s$};
 \filldraw[fill=blue!8,draw=blue!55!black]
   (0.4,0.2) rectangle (2.1,1.7);
 \node[align=center,font=\footnotesize] at (1.25,0.95) {buffer $1$\\stage $1$\\cuts};
 \filldraw[fill=blue!8,draw=blue!55!black]
   (2.5,0.2) rectangle (4.2,1.7);
 \node[align=center,font=\footnotesize] at (3.35,0.95) {buffer $0$\\stage $0$\\cuts};
 \filldraw[fill=green!8,draw=green!35!black]
   (5.05,0.2) rectangle (8.05,1.7);
 \node[align=center,font=\footnotesize] at (6.55,0.95)
   {inner controls\\$B_0,\ldots,B_l$\\main cuts};
 \filldraw[fill=orange!10,draw=orange!65!black]
   (8.5,0.2) rectangle (10.65,1.7);
 \node[align=center,font=\footnotesize] at (9.575,0.95) {end controls\\and seed};
 \draw[densely dashed,thick] (10.85,0.95) -- (12.8,0.95);
 \draw[gray] (11.85,1.78) -- (11.85,1.08);
 \node[align=center,font=\scriptsize] at (11.85,2.3)
   {arbitrarily long\\periodic continuation};
 \draw (13.05,0.12) -- (13.05,-0.12)
   node[below=3pt] {$q$};
 \draw[thick,orange!65!black] (13.05,0.25) -- (14.45,0.25);
 \node[font=\scriptsize] at (13.75,0.58) {fixed tail};
 \draw (14.45,0.12) -- (14.45,-0.12)
   node[below=3pt,align=center] {$q+b_{\rm tail}$};
 \node[anchor=west,font=\scriptsize] at (0.4,-1.0)
   {Long-continuation case shown; all solid control regions have fixed finite span.};
\end{tikzpicture}

 \caption{The order of the computation regions in one episode, not to
 scale. A marker-absent episode with a long periodic continuation is shown.
 The control regions include bounded windows for translated hypotheses.
 With an earlier first mismatch, end-control windows may extend into the
 fixed tail; all remain before the episode endpoint $q+b_{\rm tail}$.}
 \label{fig:episode-schedule}
\end{figure}

\begin{definition}
Let $E$ consist of the nonempty words starting at $s$ for which the
required data are available, the anchor rule defines $q$, and the word
ends exactly at $q+b_{\rm tail}$.  The rule is interpreted using the
placement just fixed.  Since all tests translate with the letters, it
does not depend on an ambient coordinate origin.
\end{definition}

\begin{lemma}[Prefix code]
\label{lem:prefix-code}
The episode language $E$ is a prefix code: no member is a proper prefix
of another member.  All marker and bounded control decisions associated
with an episode are unchanged when letters are appended.
\end{lemma}

\begin{proof}
By Lemma~\ref{lem:placement}, each decision window is internal before
the prescribed endpoint.  Its letters therefore remain unchanged under
extension.  A present-marker anchor is unchanged.  In the absent case,
the seed and its periodic extension are unchanged, and the first mismatch
already seen remains the first mismatch.  Thus any extension has the same
prescribed endpoint.  It can be a complete episode only if it ends there.
\end{proof}

\subsection{A guard that forces the correct end}

An \emph{exact-end test at an origin} applies the anchor rule at that
origin and requires the entire test argument to end at
$q+b_{\rm tail}$.  In particular it includes the mismatch letter when
the marker is absent.  A suffix may make this test at a hypothesized
origin preceding its own start: it must still possess all the actual
seed and inspection windows required by that hypothesis.

For a fixed finite nonempty list of candidate origins, define its
\emph{end guard} to require equal values of $\eta$ at all candidates
and to apply the exact-end test at one fixed reference candidate.  The
list and reference are fixed by the suffix test's finite choices.  All
candidate windows are mandatory under the convention above.

\begin{lemma}[Exact end guard]
\label{lem:end-guard}
Suppose a finite word begins with a complete episode $e$ of true origin
$t$ and anchor $q$.  A prefix cut $k$ lies inside this first episode.
Let a suffix argument starting at $k$ use an end guard whose candidate
list contains $t$ and has diameter at most $w$.  If the guard accepts,
the suffix ends exactly at the end of $e$, and its reported anchor is
$q$.  This conclusion holds even if the proposed suffix initially ends
before $e$ or extends past $e$ into later letters.

On the actual suffix of $e$, the guard accepts whenever $A_1(t)$ is
false, provided the placement conditions for its list hold.
\end{lemma}

\begin{proof}
\emph{Agreement of the actual anchors.}
An accepted argument contains every mandatory window, so all the
candidate marker decisions agree with the corresponding decisions in the
containing word.  The true candidate therefore has its true value of
$\eta$.  If the common value is a present marker, the reference's
anchor is immediately $q$, and the exact-end test forces the claimed end.

Suppose instead that the common value is $\bot$.  Every candidate seed
has a period less than $d_e$.  Their starts differ by at most $w$, and
\[
 K_e-w>d_e^2.
\]
Their seeds consequently determine the same two-sided periodic word by
Lemma~\ref{lem:periodic-extension}.  The overlapping union of the seeds
consists entirely of letters of this word.  The true first mismatch is after the true seed. Every later candidate
seed starts before the true seed ends, since its displacement is at most
$w<K_e$. The mismatch cannot precede the start of such a seed, and it
cannot lie in any later candidate seed: that seed consists of actual letters of the same
periodic word, whereas the mismatch letter differs from it. A candidate
seed starting earlier also ends before the true seed ends. Thus the true
first mismatch follows the entire seed union.  It is therefore also the first mismatch after the reference
seed.

\emph{Rejection of premature and excessive consumption.}
For a proposed suffix ending too soon, either a mandatory window is
absent, or no earlier reference mismatch exists to certify its proposed
endpoint.  For a suffix extending too far, the actual first mismatch is
still the first one, so its exact-end test still requires the original
endpoint $q+b_{\rm tail}$.  Both cases are excluded.

\emph{Acceptance on a stable true episode.}
Finally, if $A_1(t)$ is false, every candidate is within distance $w$ of
$t$ and has the same $\eta$ value.  On the true suffix all windows are
available by placement.  The argument above shows that the reference
exact-end rule supplies the true endpoint, so the guard accepts.
\end{proof}

Only the reference candidate makes an unbounded continuation search in
this guard.  The other candidates supply bounded marker and seed data.
At the last residual stage we will instead decode the origin from bounded
data before making its one exact-end test.  These distinctions will also
be used to compile the suffix tests at height one.

\section{Three split constructions}
\label{sec:splits}

All timing data are now fixed. We use them to define three deterministic
episode updates and their prefix/suffix tests. The main update has the
true monoid action as its linear part. Two larger affine updates encode
the main update and then each other on distinguished basis vectors.
Their graph languages are constructed in reverse order, starting from
an empty skipped class. This section proves semantic correctness of those
tests; Section~\ref{sec:compilation} establishes their expression heights.

There are three different orders to keep separate. Along the word, buffer
$1$ precedes buffer $0$ and the main segment. The total update tables are
defined in the order $g_0,g_1,g_2$. The graph tests are constructed by
stage~1, stage~0, and the main split. Every episode update is fixed on the
entire episode before a split is selected.

Throughout this section, an unavailable inspection window or an unordered
or unavailable product interval makes a test fail.  In particular, every
window belonging to a candidate list is mandatory: a test never removes
an unavailable candidate and then checks uniqueness among those remaining.
The placement in Lemma~\ref{lem:placement} supplies these windows on all
the true uses below, including the windows at false candidate origins.

\subsection{The main update and split}

Fix an episode $e\in E$, with start $s$, origin $t=s+L$, anchor $q$, and
end $b=q+b_{\rm tail}$.  Write
\[
 T_e=T_{s,b},\qquad d_j=T_{B_{j-1},B_j}\quad(1\leq j\leq l),
 \qquad \mathbf d=(d_1,\ldots,d_l).
\]
Recall that, for a residue $\ell\in\mathbb Z/n\mathbb Z$,
\[
 \mathcal E(\ell)
 =\{(\alpha,\lambda):\ell\in I_\alpha+J_\lambda\}
\]
is its clock edge set.  By Lemma~\ref{lem:clock}, this set is either
contained in the diagonal, possibly empty, or consists of one
off-diagonal edge.  With $\ell=q-t\pmod n$, define $g_{0,e}:V\to V$ by
\begin{equation}
\label{eq:main-update}
 g_{0,e}(x)=
 \begin{cases}
 xT_e+\beta(\mathbf d)T_{B_l,b},
   &\mathcal E(\ell)\text{ is contained in the diagonal},\\[2mm]
 xT_e+Y(\lambda)-X(\alpha)T_e,
   &\mathcal E(\ell)=\{(\alpha,\lambda)\},\quad\alpha\ne\lambda.
 \end{cases}
\end{equation}
Here $X(\alpha)$ and $Y(\alpha)$ denote the corresponding fields of a
tag.  In both cases this is an affine map with linear part $x\mapsto xT_e$.
All terms in its definition are fixed by the episode.

We define languages $P_x^{\rm main},Q_y^{\rm main}$ for $x,y\in V$.
In specifying a prefix test, $s$ is the start of its argument; in specifying
a suffix test, $k$ is the start of its argument.  These coordinates are
relative to the argument and do not require an external position origin.

The prefix test $P_x^{\rm main}$ chooses a tag
\[
 \alpha=(j,r,v,d_1,\ldots,d_{j-1},X,Y)
\]
and an integer $0\leq a<n$.  Its argument must end at
\[
 k=t+c_j+\delta_r a,\qquad t=s+L.
\]
It requires
\[
 k-t\pmod n\in I_\alpha,\qquad x=X,\qquad xT_{s,B_0}=v,
\]
and checks the claimed products $d_f=T_{B_{f-1},B_f}$ for $f<j$.
Only boundaries and inspections preceding its cut are used.

The suffix test $Q_y^{\rm main}$ chooses its own tag $\lambda$,
independently of the prefix choice.  First it applies the end guard of
Lemma~\ref{lem:end-guard} to the entire list
\[
 \{k-h:h\in\mathcal H\}.
\]
Thus it requires a common value of $\eta$ at these origins and applies
the exact end rule at a fixed reference origin in the list.  Denote the
resulting anchor by $q$ and the argument end by $b=q+b_{\rm tail}$.
The test requires
\[
 q-k\pmod n\in J_\lambda,\qquad y=Y(\lambda).
\]
Write the selected tag as
$\lambda=(j,r,v,d_1,\ldots,d_{j-1},X,Y)$.
For \emph{every} origin
\[
 \tau=k-c_j-\delta_r a,\qquad 0\leq a<n,
\]
the test computes $B_j(\tau),\ldots,B_l(\tau)$ by the schedule at
$\tau$, and supplies the later products
\[
 d_f^{(\tau)}=T_{B_{f-1}(\tau),B_f(\tau)}\qquad(j<f\leq l).
\]
It checks that the following function of $c\in M$ is constant:
\begin{equation}
\label{eq:main-line-test}
 c\longmapsto
 v\bigl(d_1\cdots d_{j-1}c\,d_{j+1}^{(\tau)}\cdots d_l^{(\tau)}\bigr)
 +\beta\bigl(d_1,\ldots,d_{j-1},c,
                   d_{j+1}^{(\tau)},\ldots,d_l^{(\tau)}\bigr).
\end{equation}
If its constant value is $z_\tau$, it further requires
\[
 z_\tau T_{B_l(\tau),b}=y.
\]
Empty product lists have their usual identity interpretation.  All these
checks are made within the suffix argument.  In particular, the exact end
rule is run only at the reference origin; other origins use its fixed
anchor $q$ and bounded controls.

Two different lists occur in this definition. The end guard uses every
$h\in\mathcal H$, independently of the chosen tag, so it contains the true
origin even when the suffix chooses a wrong role. The line-constancy tests
use all $n$ origins for the suffix's own role $(j,r)$. Both lists are fixed
and retain all their mandatory windows. The first enforces a true end;
the second will use product transport for completeness.

The main split skips episodes whose clock, inner transport, or end
stability may fail. Its residual decoder will use the predicate $A_0$,
but first it must obtain a common anchor from the end guard. For that
decoder we require $A_1$ to be false at every candidate origin, uniformly
over all choices of the cut. These origins are $t+p-p'$ with
$p,p'\in\mathcal P_0$, so their displacements from $t$ form $\Delta_0$.
This requirement determines the second exceptional class. Set
\begin{equation}
\label{eq:residual-domains}
 \begin{split}
 D_0&=\{e\in E:A_0(t,q)\text{ or }A_1(t)\},\\
 D_1&=\{e\in D_0:\text{there exists }\gamma\in\Delta_0
                              \text{ with }A_1(t+\gamma)\},\\
 D_2&=\varnothing.
 \end{split}
\end{equation}
Thus $D_1$ consists of those episodes already in $D_0$ for which the
all-candidate end-stability requirement fails. Since $0\in\Delta_0$,
an episode in $D_0\setminus D_1$ satisfies $A_0(t,q)$ and has $A_1$
false at every $t+\gamma$, $\gamma\in\Delta_0$. Stage~$0$ can therefore
use a common anchor and seek the true origin through $A_0$. Stage~$1$
will instead decode the bounded end-instability predicate in the
definition of $D_1$ before testing that origin's end.

The buffers will implement affine lifts $g_1:U_1\to U_1$ and
$g_2:U_2\to U_2$, defined below. The three planned split applications
are
\begin{center}
\begin{tabular}{@{}llll@{}}
\toprule
Split & Update tested & Episode class & Skipped class \\
\midrule
main & $g_0$ & $E$ & $D_0$ \\
stage $0$ & $g_1$ & $D_0$ & $D_1$ \\
stage $1$ & $g_2$ & $D_1$ & $D_2=\varnothing$ \\
\bottomrule
\end{tabular}
\end{center}
These applications are used in the order stage~1, stage~0, main.

\begin{proposition}[Main split]
\label{prop:main-split}
The languages $P_x^{\rm main},Q_y^{\rm main}$ satisfy the two split
hypotheses of Lemma~\ref{lem:split} for the episode update $g_0$, with
$D=E$ and $D'=D_0$.
\end{proposition}

\begin{proof}
\emph{Soundness for arbitrary independent tags.}
Suppose that a prefix of a word in $E^*$ factors as
$P_x^{\rm main}Q_y^{\rm main}$.  The prefix length is positive and its
cut lies before the earliest possible first episode end, by
Lemma~\ref{lem:placement}.  Hence there is a first episode $e$.
If the prefix chose $k=t+h$ with $h\in\mathcal H$, then the true origin
$t$ belongs to the suffix's end-guard list.  Its diameter is
$\operatorname{diam}\mathcal H<w$.  Lemma~\ref{lem:end-guard} therefore
forces the concatenation to end exactly at the end of $e$, with its true
anchor $q$.  This conclusion applies also to a proposed factor ending
too early or extending into later episodes; it does not presume that the
suffix has already identified the episode correctly.

Let $\alpha$ and $\lambda$ be the independently chosen prefix and suffix
tags.  The two residue tests give
$(\alpha,\lambda)\in\mathcal E(q-t)$.  If this edge set is contained in
the diagonal, then $\alpha=\lambda$.  The true origin is consequently
among the candidates for the selected role $(j,r)$ in the suffix test.
At that candidate, the transmitted vector and preceding products are
correct by the prefix checks, and the later products are correct by the
suffix checks.  Evaluating the constant function
\eqref{eq:main-line-test} at the actual $d_j$ gives the true corrected
vector at $B_l$.  Its transport to $b$ proves $g_{0,e}(x)=y$.
If instead there is one off-diagonal edge, the prefix forces
$x=X(\alpha)$ and the suffix forces $y=Y(\lambda)$.  Formula
\eqref{eq:main-update} again gives $g_{0,e}(x)=y$.  Thus every accepted
split is sound, regardless of which tags were guessed.

\emph{Completeness at one usable scale.}
Let $e\in E\setminus D_0$ and fix $x\in V$.
Then $q-t\pmod n\notin W$, some scale $r$ is usable at $t$, and
$A_1(t)$ fails.  Put $v=xT_{s,B_0}$.  Lemma~\ref{lem:affine-line}
chooses an interval $j$ for which the corrected segment function is
constant on the coordinate line through the true tuple $\mathbf d$.
Choose the tag with this $j,r,v$, the true preceding products, and
\[
 X=x,\qquad Y=g_{0,e}(x).
\]
The full diagonal is available outside $W$.  Since
$\gcd(\delta_r,n)=1$, the cuts $t+c_j+\delta_r a$ for $0\leq a<n$
realize every residue modulo $n$.  One therefore satisfies both clock
tests for this tag.

The end guard succeeds because $A_1(t)$ fails.  Each candidate for the
selected role differs from $t$ by a multiple of $\delta_r$ of magnitude
at most $R_r$.  Lemma~\ref{lem:boundary-transport} leaves all its later
interval products and its $B_l$-to-$b$ product unchanged.  Thus every
candidate in the universal suffix check has the same constant function
and the same final output as the true candidate.  Notice that this uses
only the later products: no claim about the vector at a shifted $B_0$
is needed.  The selected cut gives a split with output $g_{0,e}(x)$.
\end{proof}

\subsection{Total tables and affine lifts}

To handle $D_0$, we encode the entire main update, including its affine
correction, in the linear part of an update on $U_1$. Repeating this
construction on $U_2$ will handle the second exceptional class. The
notation $[v,c]$ always denotes a basis vector of
$U_{i+1}=\F_2^{U_i\times M}$, even when $v=0$; in particular,
$[0,1_M]$ is distinct from the zero vector of $U_{i+1}$.

For each episode, a \emph{total prefix table} at $a_i$ is a family of
maps
\[
 \Phi_i(c):U_i\longrightarrow U_i\qquad(c\in M)
\]
determined by the anchored schedule and the suffix data from $a_i$
onward, such that
$\Phi_i(T_{s,a_i})=g_{i,e}$.
The table must be defined at every $c\in M$, including values that are
not attainable by a prefix of the relevant length.

Here is the initial table explicitly.  Set
\[
 A=T_{a_0,b},\qquad C=T_{B_l,b}.
\]
The main tuple $\mathbf d$ and the clock edge set are determined from
the suffix after $a_0$.  For every $c\in M$ and $v\in U_0=V$, define
\begin{equation}
\label{eq:initial-prefix-table}
 \Phi_0(c)(v)=
 \begin{cases}
 v(cA)+\beta(\mathbf d)C,
       &\mathcal E(q-t)\text{ is contained in the diagonal},\\[2mm]
 v(cA)+Y(\lambda)-X(\alpha)(cA),
       &\mathcal E(q-t)=\{(\alpha,\lambda)\},\quad\alpha\ne\lambda.
 \end{cases}
\end{equation}
All products retain their reading order.  Substitution of
$c=T_{s,a_0}$ gives \eqref{eq:main-update}, while the formula itself
uses no knowledge of that true prefix product.

\begin{lemma}[Total tables and affine lifts]
\label{lem:table-lift}
There are episode updates $g_{1,e}:U_1\to U_1$ and
$g_{2,e}:U_2\to U_2$, affine for every $e\in E$, with the following
properties.  For $i=0,1$, the update $g_{i,e}$ has a total prefix table
$\Phi_i$ at $a_i$.  Define a linear map $H_i:U_{i+1}\to U_{i+1}$ on
basis vectors by
\[
 [v,c]H_i=[\Phi_i(c)(v),1_M].
\]
Writing
$\mathbf d_i=(T_{B_{i,0},B_{i,1}},\ldots,
               T_{B_{i,l_i-1},B_{i,l_i}})$,
the next update is
\begin{equation}
\label{eq:lift-update}
 g_{i+1,e}(z)=
 \left(z\rho_i(T_{s,a_i})
       +\beta_i(\mathbf d_i)\rho_i(T_{B_{i,l_i},a_i})\right)H_i.
\end{equation}
Its linear part $L_{i+1,e}$ satisfies
\begin{equation}
\label{eq:lift-linear-part}
 L_{i+1,e}=\rho_i(T_{s,a_i})H_i,
 \qquad
 [v,1_M]L_{i+1,e}=[g_{i,e}(v),1_M]
 \quad(v\in U_i).
\end{equation}
\end{lemma}

\begin{proof}
The initial total table is \eqref{eq:initial-prefix-table}.
Given any total table $\Phi_i$, prescribing $H_i$ on the basis defines
a unique linear map.  The maps $\Phi_i(c)$ need not be linear as maps
on $U_i$: their values only select the images of basis vectors in the
larger space.  Formula~\eqref{eq:lift-update} is therefore affine, with
the linear part displayed in \eqref{eq:lift-linear-part}.  Moreover,
\[
 [v,1_M]\rho_i(T_{s,a_i})H_i
   =[v,T_{s,a_i}]H_i
   =[\Phi_i(T_{s,a_i})(v),1_M]
   =[g_{i,e}(v),1_M].
\]

The lift has been constructed from a total table. We must still verify
that the next suffix can compute such a table without the earlier prefix.
This is the precise role of the reversed physical buffer order.

It remains to obtain the table $\Phi_1$ needed to construct $g_2$.
The buffer order is $1$ then $0$, so $a_1$ precedes all the segment
data of buffer $0$.  Define, for every $c\in M$ and $z\in U_1$,
\begin{equation}
\label{eq:lift-table}
 \Phi_1(c)(z)=
 \left(z\rho_0(cT_{a_1,a_0})
       +\beta_0(\mathbf d_0)\rho_0(T_{B_{0,l_0},a_0})\right)H_0.
\end{equation}
The actual tuple $\mathbf d_0$, the displayed suffix products, and
the entire map $H_0$ are available from the suffix after $a_1$.
Indeed, $H_0$ is built from the total table $\Phi_0$, whose values
depend only on the later anchored schedule and suffix products; it
does not require the true value of $T_{s,a_0}$. More explicitly, with
$q,t,\mathbf d,A,C$ fixed by those later data, the suffix enumerates
\eqref{eq:initial-prefix-table} for every $c'\in M$ and $v\in U_0$.
Those values give every basis image of $H_0$, including images indexed by
infeasible prefix values. Thus
\eqref{eq:lift-table} is a total table with the required dependence.
Substitution of $c=T_{s,a_1}$ yields \eqref{eq:lift-update} for $g_1$.
The preceding construction then gives $H_1$ and $g_2$.
All these maps are defined from the complete episode independently
of any split or skipped class.
\end{proof}

\begin{lemma}[Recovery on sequences]
\label{lem:sequence-lift}
For any $D\subseteq E$ and $i\in\{0,1\}$, graph tests for $g_{i+1}$
on $D^*$ yield graph tests for $g_i$ on $D^*$ by finite Boolean
operations.  Consequently this recovery does not increase generalized
star height.
\end{lemma}

\begin{proof}
For a sequence $e_1\cdots e_k\in D^*$, the linear part of the affine
composite of $g_{i+1,e_1},\ldots,g_{i+1,e_k}$ is
$L_{i+1,e_1}\cdots L_{i+1,e_k}$, in reading order.
Equation~\eqref{eq:lift-linear-part} shows that the distinguished basis
vectors return to distinguished basis vectors after every episode.
By induction, this product sends $[v,1_M]$ to
\[
 [g_{i,e_k}(\cdots g_{i,e_1}(v)\cdots),1_M].
\]
If $R_{\xi,\zeta}$ are graph tests for the affine composite on
$U_{i+1}$, the graph test for sending $v\in U_i$ to $w\in U_i$ is
therefore
\[
 \bigcup_{f\in U_{i+1}}
 \left(R_{[v,1_M],[w,1_M]+f}\cap R_{0,f}\right),
\]
by Lemma~\ref{lem:affine-recovery}.  The zero in $R_{0,f}$ is the
zero vector of $U_{i+1}$.  This formula applies to the entire sequence,
including the empty sequence, and proves the claim.
\end{proof}

\subsection{Two aligned decoders}

We have obtained affine updates and a recovery identity for entire
sequences. It remains to give the two residual splits. Their algebraic
transmission rule is the same; their origin decoders use different
predicates and a different order of end detection.

For $i=0,1$, we now define $P_x^{(i)},Q_y^{(i)}$ for inputs and outputs
$x,y\in U_{i+1}$.  These tests implement $g_{i+1}$ on $D_i$, with
skipped class $D_{i+1}$.

The prefix test chooses an offset $p\in\mathcal P_i$ and ends at
$k=s+p$.  If the kind assigned to $p$ is
\[
 (j,v,d_1,\ldots,d_{j-1}),
\]
it checks
\[
 v=x\rho_i(T_{s,B_{i,0}}),\qquad
 d_f=T_{B_{i,f-1},B_{i,f}}\quad(1\leq f<j).
\]
The relevant boundaries are the fixed buffer positions relative to
$s$, and all these checks precede the cut.

The suffix lists the candidate origins
\[
 \tau_{p'}=k+L-p'\qquad(p'\in\mathcal P_i).
\]
Its decoding and end checks are as follows.
\begin{enumerate}[label=\textup{(\arabic*)},leftmargin=*]
\item At stage $0$, require $A_1(\tau_{p'})$ to be false for every
candidate.  Apply the end guard to the full candidate list and obtain
its anchor $q$.  Require a unique $p'$ such that
$A_0(\tau_{p'},q)$ holds.  All evaluations of $A_0$ use this same
anchor; they make no additional unbounded end searches.
\item At stage $1$, first require a unique $p'$ such that
\[
 \text{there exists }\gamma\in\Delta_0
                  \text{ with }A_1(\tau_{p'}+\gamma).
\]
This test uses bounded controls and does not involve $q$.
After decoding $p'$, apply the exact end rule at $\tau_{p'}$ to
obtain $q$ and the argument end.
\end{enumerate}
At stage $0$ the candidate list has diameter $h_0<w$, as required
by the end guard.  Stage $1$ does not use that guard and needs no
such bound on $h_1$.

After decoding, set $s'=k-p'$ and use the kind attached to $p'$,
say $(j,v,d_1,\ldots,d_{j-1})$.  All subsequent buffer positions
are scheduled relative to $s'$, and the late controls use
$t'=s'+L$ and the anchor just obtained.  The suffix supplies the
later products $d_f$ for $j<f\leq l_i$ and checks that
\[
 c\longmapsto
 v\rho_i(d_1\cdots d_{j-1}c\,d_{j+1}\cdots d_{l_i})
 +\beta_i(d_1,\ldots,d_{j-1},c,d_{j+1},\ldots,d_{l_i})
 \quad(c\in M)
\]
is constant.  If its value is $z$, the test requires
\[
 \bigl(z\rho_i(T_{B_{i,l_i},a_i})\bigr)H_i=y.
\]
The map $H_i$ is computed from the explicit total tables
\eqref{eq:initial-prefix-table} and \eqref{eq:lift-table}, keeping
the decoded time and anchor data fixed.  These tables call no graph
test on an episode sequence.  The suffix uses only its own data;
if an incorrect hypothesis places required data outside its argument,
the test rejects that argument.

\begin{proposition}[Aligned splits]
\label{prop:aligned-splits}
For $i=0,1$, the languages $P_x^{(i)},Q_y^{(i)}$ satisfy the two
hypotheses of Lemma~\ref{lem:split} for the update $g_{i+1}$, with
$D=D_i$ and $D'=D_{i+1}$.
\end{proposition}

\begin{proof}
\emph{Soundness: decode the true offset and force the true end.}
Consider an accepted concatenation
$P_x^{(i)}Q_y^{(i)}$ at the start of a word in $D_i^*$.
The positive prefix cut lies before every possible first episode
end, so there is a first episode $e\in D_i$.  If the prefix chose
$p$, then its true origin occurs in the suffix list as
\[
 t=s+L=k+L-p=\tau_p.
\]

At stage $0$, the end guard therefore forces exact consumption of
this first episode and supplies its true $q$.  The suffix also
requires $A_1(t)$ to fail.  Since $e\in D_0$, this implies
$A_0(t,q)$.  The true candidate $p$ thus satisfies the decoding
predicate, so uniqueness forces the decoded $p'$ to equal $p$.

At stage $1$, membership of $e$ in $D_1$ directly implies that
$p$ satisfies the decoding predicate.  Every window in that
predicate is mandatory and, on the true episode, lies inside the
episode.  An accepted proposed suffix therefore sees the same
bounded data even if its proposed end was too early or too late.
Uniqueness first forces $p'=p$; the exact end rule at this correctly
decoded origin then forces exact consumption of $e$.

In both stages the prefix and suffix have consequently used the
same kind and the true buffer schedule.  The prefix supplies the
correct vector at $B_{i,0}$ and the preceding interval products;
the suffix supplies the correct later products.  Line constancy
then gives the true corrected vector at $B_{i,l_i}$.  Its subsequent
work action and terminal map $H_i$ give exactly
\eqref{eq:lift-update}, proving $g_{i+1,e}(x)=y$.

\emph{Completeness: count all lost offsets of the needed kind.}
Fix $e\in D_i\setminus D_{i+1}$ and
$x\in U_{i+1}$.  Apply Lemma~\ref{lem:affine-line} to the true
buffer tuple and the vector
\[
 v=x\rho_i(T_{s,B_{i,0}}).
\]
It gives an interval $j$ and thus one required kind
$(j,v,d_1,\ldots,d_{j-1})$.  Every offset of that kind passes the
prefix checks.  There are more than $C_i$ such offsets.

At stage $0$, the condition $e\notin D_1$ makes $A_1$ false at
all $t+\gamma$, $\gamma\in\Delta_0$.  For every cut $p$, all
candidate origins are of this form, so the absence checks and
the stable end guard pass and give the true $q$.  Moreover
$A_0(t,q)$ holds, so the true candidate passes the current
decoding predicate.  If an offset $p$ fails uniqueness, some
$p'\ne p$ satisfies $A_0(t+p-p',q)$.  By
Proposition~\ref{prop:inner-ambiguity}, at most $C_0$ ordered pairs
$(p,p')$ have this property.  In particular at most $C_0$ offsets
$p$ can fail uniqueness.

At stage $1$, true membership in $D_1$ supplies the true candidate
for every cut.  A failure of uniqueness gives an ordered pair
$p\ne p'$ with
\[
 \text{there exists }\gamma\in\Delta_0
             \text{ with }A_1(t+p-p'+\gamma).
\]
Lemma~\ref{lem:end-ambiguity} bounds the number of these pairs,
and hence the number of lost offsets, by $C_1$.

In either stage, a copy of the required kind remains.  At its
correctly decoded origin all later data are available by placement,
and the affine-line check gives the required output
$g_{i+1,e}(x)$.  This proves completeness.
\end{proof}

\subsection{The recovery chain}

The input and output of each split are now established. Figure~\ref{fig:proof-stages}
shows how the graph languages are assembled. A Boolean recovery changes
which update is being tested while leaving the episode domain and the
height bound fixed. A split enlarges that domain and is the only operation
in this chain that introduces an additional star.

\begin{figure}[tbp]
\centering
\begin{tikzpicture}[font=\small,x=1cm,y=1cm,
 box/.style={draw=blue!55!black,fill=blue!5,rounded corners,
 align=center,minimum width=2.6cm,minimum height=0.85cm},
 >=Stealth]
\node[box] (empty) at (0,0) {$g_2$ on $D_2^*$\\height $0$};
\node[box] (g2) at (4,0) {$g_2$ on $D_1^*$\\height $2$};
\node[box] (g1d1) at (4,-1.8) {$g_1$ on $D_1^*$\\height $2$};
\node[box] (g1d0) at (8,-1.8) {$g_1$ on $D_0^*$\\height $3$};
\node[box] (g0d0) at (8,-3.6) {$g_0$ on $D_0^*$\\height $3$};
\node[box] (g0) at (12,-3.6) {$g_0$ on $E^*$\\height $4$};
\draw[->] (empty) -- node[above,align=center,font=\scriptsize]{stage 1\\split} (g2);
\draw[->] (g2) -- node[right,font=\footnotesize]{recover $g_1$} (g1d1);
\draw[->] (g1d1) -- node[above,align=center,font=\scriptsize]{stage 0\\split} (g1d0);
\draw[->] (g1d0) -- node[right,font=\footnotesize]{recover $g_0$} (g0d0);
\draw[->] (g0d0) -- node[above,align=center,font=\scriptsize]{main\\split} (g0);
\end{tikzpicture}

\caption{Construction of the sequence graph languages. Horizontal arrows
apply a split; vertical arrows recover the lower update by finite Boolean
operations. The labels $0,2,3,4$ are the height bounds proved in
Section~\ref{sec:compilation}. Each recovery uses the same word in both
of its affine tests.}
\label{fig:proof-stages}
\end{figure}

\begin{proposition}
\label{prop:recursive-recovery}
Three applications of Lemma~\ref{lem:split}, with intervening finite
Boolean operations, construct tests for the $T$-value on $E^*$ from
the component languages defined above.
\end{proposition}

\begin{proof}
On $D_2^*=\{\eps\}$, the graph tests for $g_2$ are $\{\eps\}$
for equal input and output and empty otherwise.  Apply
Proposition~\ref{prop:aligned-splits} at stage $1$ and
Lemma~\ref{lem:split} to obtain graph tests for $g_2$ on $D_1^*$.
Lemma~\ref{lem:sequence-lift} converts these to graph tests for
$g_1$ on $D_1^*$.

Use those tests as the skipped-class tests in the stage-$0$
application of Proposition~\ref{prop:aligned-splits}.  The split
lemma gives graph tests for $g_1$ on $D_0^*$, and
Lemma~\ref{lem:sequence-lift} recovers graph tests for $g_0$ on
$D_0^*$.  These are the skipped-class tests required by
Proposition~\ref{prop:main-split}, whose application yields graph
tests for $g_0$ on $E^*$.

Finally, Lemma~\ref{lem:affine-recovery} recovers the product of the
linear parts of the maps $g_{0,e}$.  By \eqref{eq:main-update},
these are the true $T_e$ actions on $V$.  Applying their product to
the basis vector indexed by $1_M$ identifies the $T$-value of the
whole episode sequence.  Only the three stated split applications
introduce stars through the split formula; all recovery operations
are finite Boolean operations.  Section~\ref{sec:compilation}
will establish the height bounds for the component languages and
complete the count.
\end{proof}

\section{Compilation over the original alphabet}
\label{sec:compilation}

The semantic construction is complete: three splits and Boolean
recoveries determine the monoid value on episode sequences. To turn this
into a height-four expression, we must account for the component languages
and the final remainder. This section proves both claims directly over
$\Sigma$. In particular, a suffix language must have its own height-one
expression even when its guessed origin is wrong and no matching prefix
exists.

The common mechanism is a finite collection of literal word templates
$a b^h c$. Every component has at most one unbounded continuation; it
follows one fixed period until its first mismatch. On a fixed template,
all other tests become ultimately periodic conditions on the exponent
$h$. This provides the expression itself, rather than inferring a height
bound from mere regularity.

\subsection{A template lemma}

\begin{lemma}[Periodic templates]
\label{lem:templates}
Consider a finite collection of templates
\begin{equation}
 \label{eq:periodic-template}
 a b^h c,\qquad h\ge0,
\end{equation}
where $a,b,c\in\Sigma^*$ are fixed for each template and $b\ne\eps$.
Suppose a test on each template uses only the following data:
\begin{enumerate}[label=\textup{(\roman*)}]
\item letters in fixed finite collections of windows at offsets from
the argument's start or end, with mandatory availability checks;
\item order and availability of endpoints chosen from fixed finite sets
of offsets from the start or end;
\item $M$-products between these endpoints, whenever available and
ordered, and lengths modulo fixed positive integers;
\item fixed finite combinations, choices, and calculations from these
data.
\end{enumerate}
The accepted words have a generalized expression of height at most one.
Adding or removing a finite language preserves this bound.

Moreover, any of the exact-end tests of Section~\ref{sec:episodes},
with its reference or decoded origin chosen from finitely many offsets
relative to the test start, is covered by a finite language and finitely
many templates of the form~\eqref{eq:periodic-template}.
\end{lemma}

\begin{proof}
\emph{Ultimately periodic data and explicit expressions.}
Fix a template.  For sufficiently large $h$, all the indicated
availability and order relations stabilize.  Letters at a bounded
distance from the start or end also stabilize, since increasing $h$
inserts a full copy of $b$.  A product between boundedly separated
endpoints then stabilizes.  A product spanning the growing middle has
the form
\[
 u\,T(b)^{h-j}\,v
\]
for fixed $u,v\in M$ and a fixed integer $j$, once $h$ is sufficiently
large.  This follows by separating its fixed initial and terminal pieces
from its complete copies of $b$.  Powers of an element of a finite monoid
are ultimately periodic: if two powers coincide, multiplying by further
powers gives the eventual period.  Lengths modulo any fixed integer are
periodic in $h$.  Thus the entire finite data array is ultimately
periodic.  A fixed finite computation from this array, including a
universal or uniqueness check on a fixed finite list, remains ultimately
periodic.

Let the resulting accepted exponents have period $f>0$ after a threshold.
They are a finite set together with finitely many progressions
$h_0+f\mathbb N$.  Each progression contributes the expression
\[
 a\,b^{h_0}\,(b^f)^*\,c,
\]
where the displayed fixed powers are finite words.  These expressions
have height one; the exceptional words have height zero.  A finite union
proves the first assertion, including a finite number of templates and
finite exceptions.

\emph{The template cover for an independent exact-end test.}
For the final assertion, fix one possible end-rule origin relative to the
argument start.  The present-marker branch has a bounded anchor and
therefore bounded total length.  For the absent-marker branch, the seed
is at a fixed offset and has fixed length.  Make finite choices of all
letters through the seed, of its short period and phase, and of the
fixed-length tail beginning at the first mismatch.  These choices fix
a prefix $a$ through the seed and a nonempty period word $b$ beginning
at the next position.  If the mismatch occurs after $h$ full copies of
$b$ and a partial copy of length $r<|b|$, place that partial copy, the
mismatching letter, and its remaining $b_{\rm tail}-1$ letters in $c$.
There are finitely many such choices, and the mismatching letter is
required to differ from the next periodic letter.  They give exactly
templates~\eqref{eq:periodic-template} for this continuation and end.
The bounded marker and availability conditions can still be tested on
the templates.  Taking the finite union over possible end-rule origins
proves the assertion.
\end{proof}

This proof applies to the exact-end rule of a suffix considered by itself.
It does not presuppose that its guessed origin is correct, or that it is
actually the suffix of an episode.  If a hypothesized origin lacks a
mandatory seed or window, it is rejected.  Otherwise its own exact-end
rule supplies the template.

\begin{proposition}[Component bounds]
\label{prop:components}
The episode languages $E,D_0,D_1$ and every suffix component $Q_y$ in
Section~\ref{sec:splits} have generalized star height at most one.
Every prefix component $P_x$ has height zero.  These bounds hold over
the original alphabet $\Sigma$.
\end{proposition}

\begin{proof}
We first identify the single end rule that supplies each template cover,
then verify that every remaining operation is a finite calculation from
the allowed data. No part of this argument assumes that a standalone
suffix has guessed the true episode origin.

All prefix cuts come from fixed finite sets of offsets, so the lengths
of $P_x$ words are bounded.  Hence $P_x$ is finite.

Every other language in the assertion has one exact-end rule on its
entire argument.  For $E,D_0,D_1$ this is the true episode rule.  A main
or stage-0 suffix uses the reference rule in its end guard.  A stage-1
suffix uses the rule at its uniquely decoded candidate. For a stage-1
suffix considered alone, first make a finite case distinction for the
decoded offset $p'$. The uniqueness predicate inspects all candidates'
bounded end-search windows and their finitely many $\gamma$ translates;
it contains no mismatch search. The case then has just the exact-end
rule at $k+L-p'$, whether or not that is a genuine episode origin.
There are only
finitely many reference or decoded offsets; we may take a finite union
over them.  Lemma~\ref{lem:templates} therefore supplies the finite
template cover even for standalone suffixes with incorrect hypotheses.

We check that the remaining tests use only the finite data in that lemma.
Marker searches and their stability checks inspect bounded windows at
fixed nominal offsets.  So do usability, the phase and lattice choices,
and all additional candidate translates.  Every selected main boundary
lies within a fixed bounded neighborhood of its center.  Buffer
boundaries have fixed offsets.  Consequently all partial products in
the component computations have endpoints selected from finite sets of
start offsets, except for the argument end itself.

The anchor supplied by the exact-end rule is the argument end minus
$b_{\rm tail}$.  Thus the clock and $A_0$ use only fixed-modulus
differences between this end offset and their candidate origins.
The predicates $A_1$ and the stage-1 decoding predicate are wholly
bounded.  Equality, universal checks, and uniqueness over their finite
candidate lists are finite calculations, with the mandatory availability
convention retained for every candidate.

The affine-line checks range over finite vector spaces and the finite
monoid $M$.  Their data consist of the specified states and actual
partial products just described.  The tables $\Phi_i$ and $H_i$ are
computed by their explicit formulas using hypothetical elements of $M$
and actual suffix products and schedule data.  Their finite recursion
does not invoke a graph test on an episode sequence.  The total-table
property permits every hypothetical monoid value, so no extra prefix
existence test is being imposed.

Finally, no translated candidate makes a second unbounded continuation
search.  The guards compare bounded marker data and use just their one
reference rule.  Stage~1 makes bounded decoding checks before applying
the one decoded rule.  Accordingly all remaining acceptance conditions
are fixed finite calculations from the template data.  Applying
Lemma~\ref{lem:templates} proves the bounds.  Every resulting expression
uses finite words in $\Sigma^*$ and Boolean operations, concatenation,
and stars over $\Sigma$; no auxiliary tag has become an input letter.
\end{proof}

\subsection{The final remainder}

The episode rule does not require a word to end in a complete episode.
Define
\[
 F=\neg(E\Sigma^*),\qquad \Sigma^*=\neg0.
\]
Thus $F$ consists of words with no episode prefix and has height at most
one.  To finish, we need the same bound for its monoid-value tests.

\begin{lemma}[Remainder]
\label{lem:remainder}
For every $d\in M$, the language $F\cap T^{-1}(d)$ has height at most
one.  Every word admits a unique factorization into a sequence of
episodes followed by a word in $F$.
\end{lemma}

\begin{proof}
\emph{Product tests on the incomplete tail.}
Choose a fixed length threshold exceeding all the bounded decision
windows at a first start and all possible first-episode lengths in the
present-marker case.  A longer word in $F$ cannot have a present end
marker: that would already give an episode prefix.  Its marker-absent
seed has a fixed short-period continuation.  Either no mismatch has yet
occurred, or its first mismatch is at a position $q$ satisfying
\[
 |u|-q<b_{\rm tail}
\]
when the word $u$ starts at position zero.  Otherwise the prefix ending
at $q+b_{\rm tail}$ is a complete episode.  All the required bounded
windows for that prefix are internal by Lemma~\ref{lem:placement},
even when the mismatch itself is early.

In the no-mismatch case, finite choices of the bounded initial word,
period, phase, and final partial period give templates $a b^h c$.
In the mismatch case, append one of the finitely many possible
post-mismatch portions of length less than $b_{\rm tail}$; these again
give finitely many such templates.  Words below the threshold form a
finite language.  On every template the test $T(u)=d$ is an ultimately
periodic monoid-product test, so Lemma~\ref{lem:templates} gives a
height-one language with that value on the template cover.  Intersecting
with $F$ gives exactly $F\cap T^{-1}(d)$ and keeps the same bound.

\emph{Unique parsing.}
Repeatedly remove an episode prefix while one exists.  Every episode is
nonempty, so this procedure terminates, and the remaining word is in
$F$.  Lemma~\ref{lem:prefix-code} makes each removed prefix unique.
Any alternative $E^*F$ factorization must therefore remove the same
successive episodes; it also cannot stop earlier while an episode
prefix remains.  This proves uniqueness.
\end{proof}

\subsection{The height bound}

\begin{proof}[Proof of Theorem~\ref{thm:main}]
The component bounds needed for the three applications of
Lemma~\ref{lem:split} are supplied by
Proposition~\ref{prop:components}.  If the skip graph tests have height
at most $b$, that lemma gives the new bound
\[
 \max\{2,b+1\}.
\]
Indeed the episode-domain test $D^*$ has height at most two, and the
single star in the split formula lies over component and skip tests of
height at most $\max\{1,b\}$.

The stage-1 skipped class is $D_2=\varnothing$, whose graph tests consist
only of the empty-sequence tests and have height zero.  In the order
proved in Proposition~\ref{prop:aligned-splits} and
Proposition~\ref{prop:main-split}, the bounds are therefore
\[
 \underbrace{0}_{D_2^*}
 \ \longrightarrow\ 
 \underbrace{2}_{\text{stage 1}}
 \ \longrightarrow\ 
 \underbrace{3}_{\text{stage 0}}
 \ \longrightarrow\ 
 \underbrace{4}_{\text{main split}}.
\]
The affine linear-part recoveries and distinguished-basis queries use
only finite Boolean operations and do not raise these bounds.  The
linear part of the main episode update is the action of $T_e$ on $V$.
Thus its recovered sequence tests determine $T$ on $E^*$, since the
image of the identity basis vector identifies the monoid value.

Let $A_d=E^*\cap T^{-1}(d)$ be the resulting height-at-most-four tests,
and let $F_d=F\cap T^{-1}(d)$ be the height-at-most-one tests from
Lemma~\ref{lem:remainder}.  For a language recognized by $T$ with
accepting subset $S\subseteq M$, its expression is
\[
 \bigcup_{\substack{d,f\in M\\df\in S}} A_d F_f.
\]
The factorization lemma and the morphism identity prove that this is
exactly the desired language.  Concatenation and finite union do not
increase the maximum height, which is four.  All expressions are over
$\Sigma$, proving the theorem.
\end{proof}

\bibliographystyle{plainnat}
\bibliography{references}
\end{document}
